\documentclass[11pt, a4paper]{article}

% --- Encodage et langue ---
% Compatible pdflatex / xelatex / lualatex.
% Sous xelatex/lualatex, on charge fontspec (Unicode natif) ; sous pdflatex,
% on garde inputenc/fontenc + lmodern.
\usepackage{iftex}
\ifPDFTeX
  \usepackage[utf8]{inputenc}
  \usepackage[T1]{fontenc}
  \usepackage{lmodern}
  % Caractères Unicode rencontrés dans le mémoire mais non gérés par défaut.
  \DeclareUnicodeCharacter{2248}{\ensuremath{\approx}}
  \DeclareUnicodeCharacter{2265}{\ensuremath{\geq}}
  \DeclareUnicodeCharacter{2264}{\ensuremath{\leq}}
  \DeclareUnicodeCharacter{2260}{\ensuremath{\neq}}
  \DeclareUnicodeCharacter{2212}{\ensuremath{-}}
  \DeclareUnicodeCharacter{2192}{\ensuremath{\rightarrow}}
  \DeclareUnicodeCharacter{2190}{\ensuremath{\leftarrow}}
  \DeclareUnicodeCharacter{2191}{\ensuremath{\uparrow}}
  \DeclareUnicodeCharacter{2193}{\ensuremath{\downarrow}}
  \DeclareUnicodeCharacter{2194}{\ensuremath{\leftrightarrow}}
  \DeclareUnicodeCharacter{21D2}{\ensuremath{\Rightarrow}}
  \DeclareUnicodeCharacter{2026}{\ldots}
  \DeclareUnicodeCharacter{2075}{\ensuremath{^{5}}}
  \DeclareUnicodeCharacter{2076}{\ensuremath{^{6}}}
  \DeclareUnicodeCharacter{25BA}{\ensuremath{\blacktriangleright}}
  \DeclareUnicodeCharacter{25C4}{\ensuremath{\blacktriangleleft}}
  \DeclareUnicodeCharacter{2500}{-}
  \DeclareUnicodeCharacter{2502}{|}
  \DeclareUnicodeCharacter{250C}{+}
  \DeclareUnicodeCharacter{2510}{+}
  \DeclareUnicodeCharacter{2514}{+}
  \DeclareUnicodeCharacter{2518}{+}
  \DeclareUnicodeCharacter{2082}{\textsubscript{2}}
\else
  \usepackage{fontspec}
  \defaultfontfeatures{Ligatures=TeX}
  % « Arial » (équivalent libre métrique-compatible fourni par texlive-fonts-recommended).
  \setmainfont{TeX Gyre Heros}
  \setsansfont{TeX Gyre Heros}
  \setmonofont{TeX Gyre Cursor}
\fi
\usepackage[french]{babel}

% --- Caractères Unicode manquants (valable pdflatex ET xelatex/lualatex) ---
\usepackage{newunicodechar}
\newunicodechar{≈}{\ensuremath{\approx}}
\newunicodechar{≥}{\ensuremath{\geq}}
\newunicodechar{≤}{\ensuremath{\leq}}
\newunicodechar{≠}{\ensuremath{\neq}}
\newunicodechar{⇒}{\ensuremath{\Rightarrow}}
\newunicodechar{→}{\ensuremath{\rightarrow}}
\newunicodechar{←}{\ensuremath{\leftarrow}}
\newunicodechar{↑}{\ensuremath{\uparrow}}
\newunicodechar{↓}{\ensuremath{\downarrow}}
\newunicodechar{↔}{\ensuremath{\leftrightarrow}}
\newunicodechar{−}{\ensuremath{-}}
\newunicodechar{…}{\ldots}
\newunicodechar{⁵}{\ensuremath{^{5}}}
\newunicodechar{⁶}{\ensuremath{^{6}}}
\newunicodechar{►}{\ensuremath{\blacktriangleright}}
\newunicodechar{◄}{\ensuremath{\blacktriangleleft}}
\newunicodechar{─}{\texttt{-}}
\newunicodechar{│}{\texttt{|}}
\newunicodechar{┌}{\texttt{+}}
\newunicodechar{┐}{\texttt{+}}
\newunicodechar{└}{\texttt{+}}
\newunicodechar{┘}{\texttt{+}}
\newunicodechar{₂}{\textsubscript{2}}

% --- Mise en page ---
\usepackage[margin=2.5cm]{geometry}
\usepackage{setspace}
\setstretch{1.15}

% --- Titres : forcer \paragraph et \subparagraph à passer à la ligne ---
% Par défaut LaTeX rend \paragraph (= #### Markdown) en titre « run-in »,
% collé au début du paragraphe qui suit. On les redéfinit en blocs.
\usepackage{titlesec}
\titleformat{\paragraph}[block]
  {\normalfont\normalsize\bfseries}{\theparagraph}{1em}{}
\titlespacing*{\paragraph}
  {0pt}{2.25ex plus 1ex minus .2ex}{1ex plus .2ex}
\titleformat{\subparagraph}[block]
  {\normalfont\normalsize\bfseries\itshape}{\thesubparagraph}{1em}{}
\titlespacing*{\subparagraph}
  {0pt}{2ex plus 1ex minus .2ex}{0.8ex plus .2ex}

% --- Anti-débordement (inline code, longs identifiants, URLs) ---
% Pandoc convertit `code` en \texttt{...}. Par défaut, la fonte monospace
% interdit toute coupure : les identifiants type « recursive-1024-128 »
% débordent en marge.
% Stratégie :
%   (a) marge d'élasticité d'urgence + \sloppy,
%   (b) coupure aux tirets « - » réactivée dans la fonte monospace,
%   (c) le paquet « underscore » rend « \_ » sécable en mode texte,
%       ce qui permet le retour à la ligne après un « _ » dans les
%       identifiants type « entité_émettrice », sans couper le mot
%       au milieu d'une syllabe.
\setlength{\emergencystretch}{4em}
\usepackage[strings]{underscore}

\makeatletter
% Active la coupure aux tirets « - » dans la fonte monospace
\renewcommand\ttfamily{%
  \not@math@alphabet\ttfamily\mathtt
  \fontfamily\ttdefault\selectfont
  \hyphenchar\font=`\-\relax}
\makeatother

% --- Maths ---
\usepackage{amsmath, amssymb, amsfonts}

% --- Tableaux et figures ---
\usepackage{array}
\usepackage{calc}
\usepackage{booktabs}
\usepackage{longtable}
\usepackage{graphicx}
\makeatletter
\def\maxwidth{\ifdim\Gin@nat@width>\linewidth\linewidth\else\Gin@nat@width\fi}
\def\maxheight{\ifdim\Gin@nat@height>\textheight\textheight\else\Gin@nat@height\fi}
\makeatother
\setkeys{Gin}{width=\maxwidth,height=\maxheight,keepaspectratio}
\usepackage{float}

% --- Liens ---
\usepackage[colorlinks=true, linkcolor=blue!60!black, citecolor=blue!60!black, urlcolor=blue!60!black]{hyperref}

% --- Support pour les tables sans légende Pandoc 3.8+ ---
\newcounter{none}
\providecommand{\theHnone}{\thenone}

% --- Code ---
% Build uses pandoc --no-highlight, so we provide minimal Shaded / Highlighting
% environments ourselves (pandoc still wraps code blocks in them).
\usepackage{listings}
\usepackage{xcolor}
\usepackage{fancyvrb}
\usepackage{framed}
\definecolor{shadecolor}{RGB}{248,248,248}
\newenvironment{Shaded}{\begin{snugshade}}{\end{snugshade}}
\DefineVerbatimEnvironment{Highlighting}{Verbatim}{commandchars=\\\{\}}

% --- Notes de bas de page (footnotes Pandoc) ---
\usepackage{footnote}

% --- En-têtes / pieds de page ---
\usepackage{fancyhdr}
\pagestyle{fancy}
\setlength{\headheight}{28pt}
\fancyhf{}
\fancyhead[L]{\nouppercase{\leftmark}}
\fancyhead[R]{\thepage}
\renewcommand{\headrulewidth}{0.4pt}
% Classe « article » : \leftmark est alimenté par \sectionmark.
\renewcommand{\sectionmark}[1]{\markboth{#1}{}}

% --- Titre ---
\title{Jumeaux Numériques et Intelligence Artificielle Topologique :
Vers une prédiction instantanée de la résilience et des émissions
urbaines}
\author{Thomas Clerc}
\date{2026}

% --- Pandoc tight list ---
% Pandoc insère \tightlist au début des listes serrées (items sans ligne
% vide entre eux). La définition par défaut écrase tout espace vertical,
% ce qui rend les puces illisibles. On la neutralise pour conserver les
% espacements standards de \itemize.
\providecommand{\tightlist}{}

% --- Pandoc syntax highlighting macros ---

% --- CSL refs (Pandoc --citeproc) ---

\begin{document}

% --- Page de titre ---
\begin{titlepage}
\centering
\vspace*{2cm}
{\Huge\bfseries Jumeaux Numériques et Intelligence Artificielle
Topologique : Vers une prédiction instantanée de la résilience et des
émissions urbaines\par}
\vspace{1.5cm}
{\Large Mémoire\par}
\vspace{1cm}
{\large Thomas Clerc\par}
\vspace{0.5cm}
{\large 2026\par}
\vfill
\end{titlepage}

% --- Table des matières ---
\tableofcontents
\newpage

% --- Corps ---
\section{ÉCOLE HEXAGONE}\label{uxe9cole-hexagone}

\emph{Établissement d'enseignement supérieur privé}\\
\emph{Cursus Intelligence Artificielle (IA)}

\section{RAPPORT PROFESSIONNEL DE FIN
D'ÉTUDES}\label{rapport-professionnel-de-fin-duxe9tudes}

\subsection{Jumeaux Numériques et Intelligence Artificielle Topologique
: Vers une prédiction instantanée de la résilience et des émissions
urbaines}\label{jumeaux-numuxe9riques-et-intelligence-artificielle-topologique-vers-une-pruxe9diction-instantanuxe9e-de-la-ruxe9silience-et-des-uxe9missions-urbaines}

\subsubsection{(De la Micro-Simulation à l'Analyse
Spectrale)}\label{de-la-micro-simulation-uxe0-lanalyse-spectrale}

\textbf{Présenté par :} Thomas CLERC\\
\textbf{Année Universitaire :} 2025-2026

\textbf{Sous le tutorat de :}\\
* \textbf{M. Pierre UZARRALDE (IA)}, Professeur à l'École Hexagone

\textbf{Co-encadrant de recherche (VinUniversity) :}\\
* \textbf{Équipe de recherche Smart Mobility}, VinUniversity (Hanoï,
Vietnam)

\emph{Illustration de couverture : Jumeau Numérique Urbain et Signature
Spectrale de Graphe}

\newpage

\section{REMERCIEMENTS}\label{remerciements}

Je tiens à exprimer ma profonde gratitude à l'ensemble des personnes qui
ont contribué au succès de ce travail de recherche et à l'aboutissement
de mon cursus de fin d'études.

Tout d'abord, je remercie chaleureusement mon tuteur académique,
\textbf{M. Pierre Uzarralde}, pour son suivi rigoureux, ses conseils
méthodologiques et son exigence scientifique tout au long de la
rédaction de ce rapport professionnel.

Mes remerciements s'adressent également aux équipes de recherche en
mobilité intelligente de l'Université \textbf{VinUniversity} à Hanoï
(Vietnam), qui m'ont accueilli lors de mon séjour de recherche et m'ont
fourni les moyens matériels et de collecte de données nécessaires à
cette étude.

Je remercie particulièrement les groupes \textbf{Vingroup},
\textbf{VinFast} et \textbf{V-Green} pour l'accès aux données
opérationnelles du hub de recharge ultra-rapide de Sao Bien (Vinhomes
Ocean Park), ainsi que la compagnie \textbf{GSM (Green and Smart
Mobility)} pour les données relatives à leur flotte de taxis
électriques.

Enfin, je remercie l'administration et le corps enseignant de
l'\textbf{École Hexagone} pour la qualité de l'enseignement dispensé,
ainsi que ma famille et mes proches pour leur soutien indéfectible
durant cette année académique.

\newpage

\section{RÉSUMÉ \& MOTS-CLÉS}\label{ruxe9sumuxe9-mots-cluxe9s}

\subsubsection{Résumé}\label{ruxe9sumuxe9}

Ce mémoire de fin d'études présente un cadre méthodologique de recherche
pour l'évaluation et la prédiction de la résilience cinématique et
environnementale des réseaux de transport urbain face à la transition
vers l'électromobilité. La planification des infrastructures de recharge
à haute puissance en milieu dense se heurte à une double contrainte :
l'hybridation des flux cinématiques et énergétiques, et la lenteur
computationnelle des micro-simulateurs physiques multi-agents comme
SUMO. Ce travail de recherche s'articule autour d'un double objectif
complémentaire.

D'une part, nous développons des jumeaux numériques microscopiques de
haute-fidélité calibrés par vision par ordinateur (YOLOv8) sur l'avenue
Sao Bien à Hanoï (Vinhomes Ocean Park), documentant des comportements
stochastiques d'arrêt (\emph{Dwell Time}) et l'inversion modale de
trafic lors des vacances (« Holiday Reversal »).

D'autre part, nous proposons une méthode de rupture basée sur
l'\textbf{Intelligence Artificielle Topologique Spectrale}. En
exploitant la structure mathématique non-normale de la matrice
d'adjacence du réseau urbain (rayon spectral, constante de Kreiss,
perturbations de Kato) et l'algorithme XGBoost, notre modèle prédictif
direct estime de manière instantanée les émissions globales de \(CO_2\)
sous divers régimes météorologiques grâce à un paradigme de bases de
données d'apprentissage découplées. Les résultats démontrent une
précision de prédiction supérieure à 95 \% (RMSE de 15,5 tonnes de
\(CO_2\)) et un temps d'exécution de 0,2 seconde, éliminant les goulots
d'étranglement computationnels (surcharge RAM/SWAP) et ouvrant la voie à
des boucles d'optimisation hybrides (IA-SUMO) pour les planificateurs de
la ville intelligente.

\subsubsection{Mots-clés}\label{mots-cluxe9s}

Jumeau Numérique, Micro-simulation, SUMO, Analyse Spectrale, Graphes
Non-Normaux, Constante de Kreiss, Théorie des Perturbations de Kato,
XGBoost, Mobilité Durable, Vinhomes Ocean Park, Hanoï.

\newpage

\section{LISTE DES ILLUSTRATIONS}\label{liste-des-illustrations}

\begin{itemize}
\tightlist
\item
  \textbf{Figure 1 :} Morphologie géométrique et contraintes d'accès du
  hub de recharge de Sao Bien (12 bornes en épi à 135 degrés) et de la
  zone d'attente de taxis GSM le long de la voie de service.
  (\emph{Chapitre 8})
\item
  \textbf{Figure 2 :} Représentation schématique du flux de circulation
  et du tissage des motocyclistes au niveau de l'insertion sur l'avenue
  Sao Bien. (\emph{Chapitre 8})
\item
  \textbf{Figure 3 :} Répartition modale comparée du fluide de trafic
  lors d'un midi de semaine classique (\emph{Regular Midday Baseline})
  face à l'inversion modale observée en période de congés (\emph{Holiday
  Reversal}). (\emph{Chapitre 8})
\item
  \textbf{Figure 4 :} Pipeline d'apprentissage découplé pour la
  prédiction directe des émissions de \(CO_2\) sous divers régimes
  météorologiques régionalisés. (\emph{Chapitre 7})
\end{itemize}

\newpage

\section{LISTE DES TABLEAUX}\label{liste-des-tableaux}

\begin{itemize}
\tightlist
\item
  \textbf{Tableau 1 :} Caractéristiques topologiques générales des 21
  villes étudiées (nombre de nœuds, d'arêtes, densité de connexion,
  degré moyen, indice d'asymétrie, ratio de sources et de puits).
  (\emph{Chapitre 6})
\item
  \textbf{Tableau 2 :} Métriques spectrales de non-normalité des
  matrices d'adjacence non pondérées pour les 21 réseaux urbains
  (asymétrie, rayon spectral, valeur singulière maximale, norme \(H_2\),
  constante de Kreiss). (\emph{Chapitre 6})
\item
  \textbf{Tableau 3 :} Performances comparatives des modèles
  d'apprentissage machine pour la prédiction des émissions de \(CO_2\)
  (RMSE, MAE, \(R^2\) de XGBoost, Ridge, Random Forest et MLP).
  (\emph{Chapitre 7})
\item
  \textbf{Tableau 4 :} Profil d'importance des variables (\emph{Feature
  Importance}) dans le modèle final de prédiction de \(CO_2\) par
  XGBoost. (\emph{Chapitre 7})
\item
  \textbf{Tableau 5 :} Taille géométrique des fichiers net.xml sur
  disque et occupation correspondante de la mémoire RAM en Python
  (filtrage DOM via sumolib) pour cinq villes types. (\emph{Chapitre 9})
\item
  \textbf{Tableau 6 (Annexe A) :} Indicateurs macroscopiques des 17
  simulations complètes exécutées dans le framework multi-villes
  agnostique (véhicules, vitesse moyenne, \(CO_2\) émis, conditions
  météo, taux d'adoption EV, temps d'exécution). (\emph{Annexes, Annexe
  A})
\item
  \textbf{Tableau 7 (Annexe B) :} Répartition des tonnages de \(CO_2\)
  émis par classe de véhicule de la flotte (voitures, motos, bus,
  camions) pour les 17 simulations de référence. (\emph{Annexes, Annexe
  B})
\item
  \textbf{Tableau 8 (Annexe C) :} Profil de performance informatique
  détaillé des exécutions CPU (temps de routage Dijkstra, calcul
  physique SUMO, parsing XML, durée totale). (\emph{Annexes, Annexe C})
\end{itemize}

\newpage

\section{INTRODUCTION GÉNÉRALE}\label{introduction-guxe9nuxe9rale}

\subsubsection{Contexte global : Expansion urbaine, transition
électrique et défis de la mobilité
durable}\label{contexte-global-expansion-urbaine-transition-uxe9lectrique-et-duxe9fis-de-la-mobilituxe9-durable}

Aujourd'hui, plus de la moitié de la population mondiale réside en
milieu urbain. Selon les rapports officiels des Nations Unies
(\emph{United Nations - World Urbanization Prospects}, 2018), 55 \% de
la population globale vivait en ville en 2018, et cette proportion
devrait franchir le seuil des 68 \% d'ici 2050. Cette transition
démographique sans précédent représente un apport de plus de 2,5
milliards d'urbains supplémentaires, dont près de 90 \% de la croissance
se concentrera dans les pays en développement d'Asie et d'Afrique.

Cette métamorphose rapide des métropoles s'accompagne d'une
densification extrême de l'espace et d'une augmentation géométrique des
flux de transport et des déplacements quotidiens (World Bank, 2024). On
observe l'émergence et le développement fulgurant de nouveaux quartiers
planifiés en quelques années seulement. C'est le cas emblématique du
complexe résidentiel et commercial \emph{Vinhomes Ocean Park} (VHOP),
situé en périphérie de la capitale vietnamienne, Hanoï. Conçu sur
d'anciennes parcelles agricoles, ce quartier moderne héberge déjà près
de 60 000 résidents permanents et prévoit d'en accueillir plus de 200
000 d'ici 2030 (Vinhomes, 2024).

Les infrastructures et réseaux routiers existants n'ayant pas été conçus
pour absorber une telle cinématique de croissance, cette saturation
engendre des congestions chroniques de trafic et une hausse massive des
émissions polluantes. À l'échelle mondiale, le secteur des transports
routiers constitue l'un des principaux contributeurs au changement
climatique. Selon les données de l'Agence Internationale de l'Énergie
(\emph{International Energy Agency - IEA}, 2025), le transport routier
est responsable de plus de 6 Gt de \(CO_2\) émises en 2024, affichant
une croissance continue de 8 \% depuis 2015. Dans ce bilan carbone, les
véhicules particuliers et utilitaires légers représentent plus de 60 \%
des émissions, tandis que les véhicules lourds et camions comptent pour
environ un tiers.

La congestion routière aggrave considérablement cette situation : les
cycles d'arrêt-démarrage et l'inactivité des moteurs thermiques dans les
embouteillages (\emph{idling}) provoquent une surconsommation de
carburant et des émissions de polluants locaux majeures, en particulier
chez les poids lourds. L'évaluation de l'Organisation de Coopération et
de Développement Économiques (\emph{OECD}, 2024) et de l'Agence
Européenne pour l'Environnement (\emph{EEA}, 2024) montre que
l'efficacité écologique globale ne dépend pas uniquement de l'efficacité
énergétique individuelle des véhicules, mais également de la résilience
cinématique et de la fluidité des réseaux routiers eux-mêmes, qui
peinent à dissiper les ondes de congestion en conditions de haute
densité.

Pour endiguer ce fléau environnemental, l'électrification massive des
flottes de véhicules s'est imposée comme le pivot central des politiques
publiques. En France, l'État encourage activement cette transition par
des aides financières ciblées (bonus écologique, prime à la conversion)
(\emph{Service Public}, 2026) et un déploiement massif d'infrastructures
de recharge porté par des initiatives publiques et privées. Le réseau
national s'étend rapidement le long des grands axes autoroutiers, mais
s'implante aussi dans les zones commerciales à travers les réseaux de
grande distribution, à l'image des stations installées par des enseignes
comme Lidl (Avere-France, 2025).

Au Vietnam, dans les quartiers modèles comme Vinhomes Ocean Park, la
transition vers l'électromobilité se déploie à un rythme encore plus
avancé. VHOP intègre des flottes entièrement décarbonées de transports
publics (\emph{VinBus}, opérant des bus électriques intelligents sur
site depuis 2021) (\emph{VinBus}, 2021) et de taxis électriques de la
compagnie GSM (\emph{Green SM}, 2025). Néanmoins, cette adoption
ultra-rapide d'équipements électriques déplace le verrou. La recharge
simultanée de flottes d'autobus et de taxis au niveau de stations à
haute puissance --- telles que le hub de 12 bornes rapides (150 kW DC)
de Vincom Mega Mall à Sao Bien --- crée des perturbations majeures sur
le réseau routier local. Les manœuvres d'insertion et d'attente à
proximité des hubs de recharge modifient l'écoulement naturel du trafic,
engendrant des congestions localisées inattendues.

Dès lors, la problématique générale de ce travail de recherche s'établit
comme suit : \textbf{comment développer un modèle numérique capable
d'estimer de manière instantanée, précise et généralisable les émissions
de \(CO_2\) urbaines générées par le trafic routier, dans n'importe
quelle ville du monde, sous divers scénarios de transition de flotte et
d'infrastructures ?}

\subsubsection{Problème Scientifique et Technique : La barrière
computationnelle de la micro-simulation
physique}\label{probluxe8me-scientifique-et-technique-la-barriuxe8re-computationnelle-de-la-micro-simulation-physique}

Les réseaux routiers urbains se comportent comme des systèmes complexes
hautement non-linéaires (\emph{Complex Adaptive Systems}, \emph{Urban
Dynamics}). Une modification locale de la topologie --- qu'il s'agisse
de l'ajout d'un giratoire, du retrait d'une voie de circulation pour y
aménager une piste cyclable, ou de l'installation d'un hub de recharge
rapide --- peut déclencher des ondes de choc cinématiques et propager
des congestions à l'ensemble de la ville (Braess, 1968). Face à ces
comportements émergents, les urbanistes et planificateurs sont
historiquement confrontés à une impossibilité de prévoir a priori la
robustesse et la vulnérabilité d'un réseau viaire sous charge.

Pour étudier ces interactions physiques à l'échelle granulaire,
l'ingénierie du transport s'est structurée autour d'outils de
micro-simulation multi-agents tels que SUMO (Simulation of Urban
MObility) (Krajzewicz et al., 2012), Aimsun, MATsim ou PTV Vissim. Ces
simulateurs modélisent de manière granulaire le déplacement de chaque
agent (vitesse, changement de voie, distance de sécurité) seconde par
seconde en s'appuyant sur des lois physiques de poursuite
(car-following). Couplés à des bases de données d'émissions de référence
telles que le modèle HBEFA3 (Handbook Emission Factors for Road
Transport) (Kratzsch et al., 2020), ils offrent une fidélité remarquable
pour quantifier les surémissions de \(CO_2\) associées aux comportements
transitoires des véhicules (cycles d'arrêt-démarrage).

Néanmoins, cette précision physique se heurte à une barrière
computationnelle majeure. La résolution séquentielle des équations
différentielles cinématiques pour chaque véhicule est extrêmement
coûteuse en temps de calcul CPU et en mémoire vive RAM. Si ces
simulateurs se révèlent efficaces pour des évaluations ciblées à petite
échelle (comme l'analyse locale du hub de recharge de Sao Bien à Hanoï
ou le réseau de taille modérée de Versailles), leur passage à l'échelle
métropolitaine s'avère impossible en pratique. Simuler un scénario de
plusieurs heures impliquant des centaines de milliers d'agents sur une
métropole comme Los Angeles ou Paris sature les ressources matérielles
des ordinateurs portables classiques des planificateurs, entraînant des
temps de calcul préjudiciables (plusieurs heures par run) et des
écritures de mémoire virtuelle lentes sur le disque (mécanisme de SWAP).
Cette inertie computationnelle empêche l'exploration de larges espaces
de scénarios et interdit toute utilisation pour la prise de décision en
temps réel ou dans des boucles d'optimisation automatique.

Cette limitation majeure justifie le développement de nouvelles
approches de rupture, capables de contourner la simulation physique en
s'appuyant sur l'intelligence artificielle pour prédire de manière
instantanée la pollution urbaine globale.

\subsubsection{Double Objectif de la Recherche et Positionnement
Académique}\label{double-objectif-de-la-recherche-et-positionnement-acaduxe9mique}

Pour briser ce verrou technologique, ce mémoire propose une méthodologie
hybride alliant la précision locale de la micro-simulation et la
rapidité de prédiction instantanée de l'intelligence artificielle. Nos
travaux s'articulent autour d'un double objectif :

\begin{enumerate}
\def\labelenumi{\arabic{enumi}.}
\tightlist
\item
  \textbf{La modélisation et calibration locale par Jumeau Numérique} :
  Nous construisons un jumeau numérique microscopique haute-fidélité de
  l'avenue Sao Bien à Hanoï (Vinhomes Ocean Park), calibré par vision
  par ordinateur (YOLOv8) à partir d'observations vidéo réelles. Ce
  premier axe permet d'étudier scientifiquement le comportement
  stochastique des sessions de recharge (\emph{Stochastic Dwell Time})
  et les phénomènes d'inversion de la composition de la flotte pendant
  les congés nationaux (\emph{Holiday Reversal}).
\item
  \textbf{La généralisation et l'inférence instantanée par IA
  Topologique Spectrale} : Afin de s'affranchir du coût de calcul de
  SUMO pour l'évaluation de nouvelles villes, nous développons un
  métamodèle d'apprentissage machine (XGBoost) entraîné sur un corpus
  multi-villes à grande échelle. L'hypothèse scientifique centrale de ce
  travail stipule que la topologie structurelle d'un réseau routier
  dicte ses performances dynamiques sous charge. En extrayant les
  caractéristiques spectrales (rayon spectral, constante de Kreiss,
  perturbations de Kato, normes de Hardy) de la matrice d'adjacence
  non-normale associée au graphe de la ville (Trefethen \& Embree, 2005
  ; Kato, 1995), notre IA estime en une fraction de seconde (0,2
  seconde) les émissions de \(CO_2\) pour n'importe quelle ville cible,
  sans nécessiter de simulation physique préalable.
\end{enumerate}

\subsubsection{Annonce Détaillée de la Structure du
Mémoire}\label{annonce-duxe9tailluxe9e-de-la-structure-du-muxe9moire}

Pour exposer notre recherche avec toute la rigueur universitaire
requise, ce mémoire est structuré selon le plan logique suivant :

La \textbf{Première Partie} (Chapitres 1 à 4) est dédiée à l'acquisition
des données et à la modélisation de la transition électrique. Après
avoir présenté les contextes démographiques de la transition électrique
(Chapitre 1) et notre protocole de collecte par vision par ordinateur à
Hanoï (Chapitre 2), nous détaillons le traitement, l'audit et la
correction des données de trafic brutes (Chapitre 3). Enfin, le Chapitre
4 formalise la constitution de notre corpus d'apprentissage à grande
échelle par le biais de simulations multi-villes découplées selon les
conditions météorologiques.

La \textbf{Deuxième Partie} (Chapitres 5 à 7) présente les fondations
physiques et théoriques de nos outils. Le Chapitre 5 détaille le
fonctionnement interne du simulateur SUMO (physique comportementale de
Krauß et algorithmes de routage). Le Chapitre 6 pose le formalisme
mathématique de la topologie spectrale non-normale des graphes orientés
et la théorie des perturbations de Kato. Le Chapitre 7 présente
l'architecture de notre métamodèle IA XGBoost direct à 43 features et sa
capacité de généralisation.

La \textbf{Troisième Partie} (Chapitres 8 et 9) confronte nos théories
aux applications empiriques. Le Chapitre 8 détaille l'étude de cas
microscopique du hub de recharge de Hanoï et l'évaluation des scénarios
de mitigation. Le Chapitre 9 déploie un cadre d'expérimentation global
comparant les performances computationnelles (RAM, temps CPU, SWAP) et
la précision prédictive de notre IA face au simulateur physique SUMO sur
six villes mondiales de morphologies contrastées.

\newpage

\section{PREMIÈRE PARTIE : ACQUISITION DES DONNÉES, PROTOCOLE DE
COLLECTE ET CONTEXTES DE
MOBILITÉ}\label{premiuxe8re-partie-acquisition-des-donnuxe9es-protocole-de-collecte-et-contextes-de-mobilituxe9}

\subsection{Chapitre 1 : Contexte démographique et modélisation de la
transition
électrique}\label{chapitre-1-contexte-duxe9mographique-et-moduxe9lisation-de-la-transition-uxe9lectrique}

\subsubsection{Le paradigme des villes intelligentes et de la
décarbonation des systèmes de
transport}\label{le-paradigme-des-villes-intelligentes-et-de-la-duxe9carbonation-des-systuxe8mes-de-transport}

La transition vers la ville intelligente (``Smart City'') ne se résume
pas à l'intégration superficielle de technologies de communication au
sein de l'espace urbain ; elle constitue une réponse structurelle à
l'impératif de décarbonation. Dans l'architecture d'une métropole
moderne, les systèmes de transport représentent l'un des principaux
réservoirs d'optimisation carbone. Les politiques publiques
traditionnelles, basées sur l'expansion continue des capacités de
voirie, ont montré leurs limites en se heurtant au phénomène de demande
induite. La décarbonation requiert donc une approche combinée : la
transition technologique des moteurs thermiques vers des systèmes de
propulsion électrique, et l'optimisation dynamique des flux pour
minimiser les pertes énergétiques associées à la congestion.

Cette restructuration des systèmes de transport impose de repenser la
relation entre l'espace routier et le réseau d'alimentation en énergie.
Que ce soit à Paris ou à Hanoï, les bornes de recharge s'insèrent dans
le domaine public, sur les places de stationnement ou les voies de
desserte commerciale. Cette hybridation physique transforme chaque point
de recharge en un nœud d'attraction cinématique, modifiant localement la
trajectoire des véhicules et perturbant l'écoulement des flux
environnants.

\subsubsection{Le cadre d'étude vietnamien : Vinhomes Ocean Park (VHOP)
et sa cinétique de
croissance}\label{le-cadre-duxe9tude-vietnamien-vinhomes-ocean-park-vhop-et-sa-cinuxe9tique-de-croissance}

Le développement rapide de l'Asie du Sud-Est s'accompagne d'une
transition urbaine caractérisée par la construction de villes nouvelles
planifiées en périphérie des centres historiques. Au Vietnam, ce modèle
trouve son illustration dans le projet \emph{Vinhomes Ocean Park}
(VHOP), un complexe de 420 hectares situé à l'est de Hanoï. Conçu
initialement comme une cité satellite résidentielle pour soulager la
pression démographique du district historique de Hoan Kiem, VHOP a été
dimensionné pour accueillir une population nominale de 90 000 habitants.

Cependant, la cinétique d'occupation de VHOP a dépassé les prévisions
initiales. En 2023, le nombre de résidents permanents a franchi le seuil
des 60 000 personnes, et les projections démographiques révisées
indiquent que le site pourrait accueillir à terme plus de 200 000
résidents d'ici 2030 en raison de l'attractivité des infrastructures
intégrées. En 2023, la population active présente sur site ne
représentait que 30 \% à 33 \% de la capacité finale projetée. L'analyse
des flux actuels montre que le réseau routier interne, bien que moderne,
approche déjà de la saturation lors des heures de pointe.

\subsubsection{L'intégration des flottes électriques (VinFast) et la
problématique infrastructurelle des hubs de
recharge}\label{lintuxe9gration-des-flottes-uxe9lectriques-vinfast-et-la-probluxe9matique-infrastructurelle-des-hubs-de-recharge}

Au cœur de la stratégie de mobilité de Vinhomes Ocean Park se trouve
l'écosystème électrique porté par le groupe \emph{Vingroup}, à travers
sa filiale automobile \emph{VinFast} et sa filiale de recharge
\emph{V-Green}. Les transports en commun internes sont assurés par des
bus électriques (\emph{VinBus}), et la flotte de taxis opérant sur le
site est constituée en majorité de véhicules électriques de la compagnie
\emph{GSM (Green and Smart Mobility)}.

Le cas d'étude de Sao Bien se focalise sur la station de recharge
installée en bordure du \emph{Vincom Mega Mall}, l'un des centres
d'attraction majeurs du complexe. Ce hub est équipé de 12 bornes de
recharge ultra-rapides d'une puissance unitaire de 150 kW DC.
L'intégration d'une telle infrastructure génère des contraintes
opérationnelles. D'une part, la demande électrique cumulée de la station
peut atteindre 1,8 MW en période de pointe, imposant des contraintes de
dimensionnement sur le réseau de distribution moyenne tension local.
D'autre part, la géométrie d'accès à la station, configurée en
stationnement en épi à 135 degrés le long d'une voie de service étroite,
crée des conflits d'usage avec la file des taxis en attente de clients
pour le centre commercial.

\newpage

\subsection{Chapitre 2 : Méthodologie de collecte de données terrain par
vision par
ordinateur}\label{chapitre-2-muxe9thodologie-de-collecte-de-donnuxe9es-terrain-par-vision-par-ordinateur}

\subsubsection{L'infrastructure d'observation : Stratégie du 3ème étage
(Vincom Mega Mall) et minimisation de
l'occlusion}\label{linfrastructure-dobservation-stratuxe9gie-du-3uxe8me-uxe9tage-vincom-mega-mall-et-minimisation-de-locclusion}

La collecte de données précises sur les flux de trafic en milieu urbain
dense nécessite une méthodologie d'observation rigoureuse pour garantir
la qualité des données d'entrée du jumeau numérique (voir Figures D.1 et
D.2, Annexe D). Pour le site de Sao Bien, nous avons établi un poste
d'observation temporaire au 3ème étage du \emph{Vincom Mega Mall}, au
niveau de la zone de restauration faisant face au corridor de
circulation principal et au hub de recharge.

Ce positionnement en hauteur (angle d'observation incliné entre 30 et 45
degrés par rapport à l'horizontale) répond à une contrainte technique
majeure : \textbf{la minimisation de l'occlusion visuelle}. Dans un
contexte de trafic mixte caractéristique du Vietnam, où les flux
comportent une part importante de deux-roues motorisés circulant de
front avec des autobus et des SUV de gabarit important, une capture au
niveau du sol souffre d'un biais de masquage systématique. La
perspective plongeante du 3ème étage permet d'obtenir une vue dégagée
``pseudo-orthogonale'' du réseau, garantissant que chaque trajectoire de
véhicule reste visible tout au long du segment d'analyse.

\subsubsection{Contraintes opérationnelles de prise de vue : Choix du
format Portrait vs Paysage face aux obstacles
géométriques}\label{contraintes-opuxe9rationnelles-de-prise-de-vue-choix-du-format-portrait-vs-paysage-face-aux-obstacles-guxe9omuxe9triques}

L'établissement de la ligne de visée depuis le 3ème étage du centre
commercial a imposé des contraintes géométriques strictes liées à
l'architecture du bâtiment. La présence de piliers de soutien en béton,
de montants de fenêtres en aluminium et de vitrines de magasins
obstruait une grande partie du champ visuel horizontal.

L'analyse comparative des formats de capture a révélé les éléments
suivants : * \textbf{Le format paysage (horizontal) :} Bien qu'adapté
pour capturer la longueur du tronçon routier, il intégrait dans le cadre
plusieurs obstacles physiques majeurs qui divisaient la zone d'intérêt
en sous-sections disjointes, empêchant le suivi continu des trajectoires
par l'algorithme de vision par ordinateur. * \textbf{Le format portrait
(vertical) :} En orientant la caméra verticalement, nous avons pu
aligner le champ de vision principal dans l'espace situé entre deux
piliers consécutifs du bâtiment. Cette configuration a permis de cadrer
de manière ininterrompue les trois zones critiques du corridor : la zone
d'approche amont, la zone d'insertion du hub de recharge, et la zone de
sortie vers le carrefour aval.

\subsubsection{Le pipeline de détection automatique et de classification
de véhicules via
YOLOv8}\label{le-pipeline-de-duxe9tection-automatique-et-de-classification-de-vuxe9hicules-via-yolov8}

Pour automatiser l'extraction des données de trafic à partir des
séquences vidéo haute définition capturées sur site, nous avons déployé
un pipeline de traitement d'images basé sur le modèle de réseau de
neurones convolutifs \textbf{YOLOv8} (You Only Look Once, version 8).

Ce pipeline de vision par ordinateur fonctionne de la manière suivante :
1. \textbf{Segmentation temporelle :} Les vidéos brutes sont découpées
en segments de 10 minutes pour correspondre aux intervalles standards
d'analyse de trafic. 2. \textbf{Inférence et détection :} Le modèle
YOLOv8 traite les trames vidéo avec un seuil de confiance de détection
fixé à \(0.50\). 3. \textbf{Classification catégorielle :} Les objets
détectés sont classés selon trois classes de véhicules : \emph{Standard
Car}, \emph{Electric Vehicle}, et \emph{Motorcycle}. 4. \textbf{Suivi de
trajectoire (Tracking) :} L'algorithme associe un identifiant unique à
chaque boîte de délimitation à l'aide d'un filtre de Kalman et d'une
matrice de coût basée sur le recouvrement spatial des boîtes
(Intersection over Union - IoU) entre trames successives.

\newpage

\subsection{Chapitre 3 : Traitement, audit et correction des données de
trafic
brutes}\label{chapitre-3-traitement-audit-et-correction-des-donnuxe9es-de-trafic-brutes}

\subsubsection{Analyse des biais de classification de l'IA et
application du facteur de correction
systématique}\label{analyse-des-biais-de-classification-de-lia-et-application-du-facteur-de-correction-systuxe9matique}

Malgré l'efficacité de l'architecture YOLOv8 pour la détection
automatique, la phase d'audit qualité des données a mis en évidence un
biais de classification systématique lors de l'analyse des flux de
véhicules particuliers. En comparant les résultats de l'extraction
automatisée à un comptage manuel de référence effectué sur 3 heures de
vidéo, nous avons identifié une \textbf{surestimation constante de la
classe des voitures de l'ordre de +30 \%}.

L'analyse d'erreur a révélé que ce biais provenait de deux facteurs : *
La confusion visuelle induite par les longs véhicules (SUV familiaux
VinFast de type VF8 et VF9, mini-fourgonnettes et véhicules de livraison
légers) qui, sous certains angles de vue inclinés, étaient fragmentés
par l'algorithme en plusieurs boîtes de délimitation distinctes. *
L'occlusion partielle transitoire causée par le passage de bus
électriques, qui masquaient puis révélaient des véhicules adjacents,
entraînant une réinitialisation de l'identifiant du véhicule.

Pour stabiliser le jeu de données et éliminer ce bruit de détection,
nous avons intégré un \textbf{facteur de correction mathématique
systématique de -30 \%} appliqué à la classe des véhicules particuliers
dans le pipeline de traitement des données.

\subsubsection{Établissement des baselines de trafic et d'émissions
(Midday vs Rush
Hour)}\label{uxe9tablissement-des-baselines-de-trafic-et-duxe9missions-midday-vs-rush-hour}

L'analyse des données corrigées a permis de définir les caractéristiques
de base de la mobilité sur le corridor de Sao Bien à travers deux
périodes représentatives de la journée : le creux de la mi-journée
(Midday) et le pic de fin d'après-midi (Rush Hour).

\paragraph{Baseline Midday (12:00 PM)}\label{baseline-midday-1200-pm}

Cette période caractérise un régime de trafic régulier et modéré, avec
un débit moyen observé de \textbf{124,4 véhicules par intervalle de 10
minutes}. La composition de la flotte se répartit comme suit : *
\emph{Motorcycles} : 63,6 unités (soit 51,1 \% de la flotte). *
\emph{Electric Vehicles (EV)} : 32,5 unités (soit 26,1 \%). *
\emph{Standard Cars (ICE)} : 28,3 unités (soit 22,8 \%).

Le trafic à midi est fluide, et la vitesse moyenne des véhicules
avoisine les 35 km/h. La demande de recharge au hub Vincom Mega Mall est
stable, maintenant un taux d'occupation des bornes de l'ordre de 75 \%
(9 bornes occupées sur 12).

\paragraph{Baseline Rush Hour (5:00
PM)}\label{baseline-rush-hour-500-pm}

Le pic de fin d'après-midi se caractérise par une augmentation de
\textbf{72 \% du volume global}, atteignant un débit moyen de
\textbf{214,5 véhicules par intervalle de 10 minutes}. La structure de
la flotte se modifie significativement : * \emph{Motorcycles} : 136,4
unités (soit 63,6 \% de la flotte). * \emph{Electric Vehicles (EV)} :
45,5 unités (soit 21,2 \%). * \emph{Standard Cars (ICE)} : 32,6 unités
(soit 15,2 \%).

Le volume élevé de deux-roues motorisés (63,6 \%) sature l'espace de
circulation. Les motos occupent les espaces inter-voies, créant des
trajectoires complexes d'évitement pour les véhicules plus volumineux.

\subsubsection{Phénoménologie du trafic observée : Le phénomène de
``Holiday Reversal'' et dynamiques de congestion
locales}\label{phuxe9nomuxe9nologie-du-trafic-observuxe9e-le-phuxe9nomuxe8ne-de-holiday-reversal-et-dynamiques-de-congestion-locales}

L'analyse des données de collecte a mis en lumière un comportement de
trafic singulier lors des périodes de congés officiels (notamment les
journées du 30 avril et du 1er mai, correspondant aux vacances
nationales au Vietnam). Nous avons qualifié ce phénomène de
\textbf{``Holiday Reversal''} (inversion de vacances).

Alors que les jours de semaine classiques présentent une domination
nette des deux-roues motorisés (51 \% à 64 \% de la flotte), les jours
fériés révèlent une structure inversée, dite ``Car-Locked'' (voir Figure
D.3, Annexe D). Les observations enregistrées à 12h00 lors des vacances
indiquent : * Une baisse de la proportion de motos, qui chute à
\textbf{34,6 \%} (soit 42,25 unités). * Une augmentation importante de
la part des véhicules électriques, qui grimpe à \textbf{43,3 \%} (soit
53,00 unités). * Une part stable de voitures thermiques à \textbf{19,4
\%} (soit 23,75 unités).

Ce changement s'explique par deux facteurs d'usage : d'une part, les
résidents privilégient les déplacements en voiture familiale pour se
rendre au centre commercial ; d'autre part, la flotte de taxis
électriques GSM intensifie son activité pour répondre à la demande des
visiteurs externes (Apur, 2025). Pour le hub de recharge, cette
configuration représente une situation critique : le taux d'occupation
des bornes atteint 100 \% de manière quasi-continue.

\newpage

\subsection{Chapitre 4 : Constitution du corpus d'apprentissage par
simulation multi-villes à grande
échelle}\label{chapitre-4-constitution-du-corpus-dapprentissage-par-simulation-multi-villes-uxe0-grande-uxe9chelle}

\subsubsection{Nécessité d'un corpus de données synthétiques
multi-villes}\label{nuxe9cessituxe9-dun-corpus-de-donnuxe9es-synthuxe9tiques-multi-villes}

Bien que les données d'observation réelles recueillies à Hanoï (via
YOLOv8) apportent une vérité terrain indispensable pour la calibration
fine des comportements locaux, elles s'avèrent structurellement
insuffisantes pour entraîner un modèle d'intelligence artificielle
prédictif à vocation généralisable. Une campagne de mesure sur site, par
nature figée dans l'espace et le temps, ne permet pas d'observer la
diversité topologique indispensable (réseaux radiaux, orthogonaux,
mixtes) ni de faire varier de manière systématique les charges de trafic
et les taux catégoriels d'électrification de la flotte.

Pour s'affranchir de cette contrainte empirique et entraîner notre
modèle prédictif spectral XGBoost, nous devons constituer un corpus
d'apprentissage étendu couvrant des milliers de situations de
congestion. C'est à ce stade qu'intervient le couplage avec la
micro-simulation (SUMO), exploitée non pas comme un outil d'analyse
locale, mais comme un \textbf{générateur de données synthétiques à haute
fidélité physique} (Physical Proxy). En simulant un large éventail de
scénarios sur plusieurs villes du globe, nous créons la base de données
d'apprentissage nécessaire pour nourrir notre métamodèle d'IA.

\subsubsection{Sélection des environnements urbains et diversité
topologique}\label{suxe9lection-des-environnements-urbains-et-diversituxe9-topologique}

Afin de garantir la robustesse et la capacité de généralisation
transversale (\emph{cross-city}) de notre modèle, nous avons sélectionné
un ensemble représentatif de villes mondiales présentant des géométries
de voirie contrastées. Cette diversité topologique permet à l'IA
d'apprendre la relation entre la structure géométrique intrinsèque d'un
graphe et ses performances cinématiques :

\begin{itemize}
\tightlist
\item
  \textbf{Versailles} : Réseau de taille modérée (2 308 nœuds, 9 095
  arêtes) caractérisé par une structure planifiée régulière et aérée.
\item
  \textbf{Paris} : Réseau historique de grande taille (15 735 nœuds, 62
  378 arêtes) présentant une structure radiale dense avec de nombreux
  ronds-points complexes.
\item
  \textbf{Madrid} : Réseau métropolitain étendu (61 258 nœuds, 214 400
  arêtes) combinant un centre-ville dense et un maillage autoroutier
  périphérique.
\item
  \textbf{Berlin} : Réseau mixte de grande dimension (61 135 nœuds, 213
  900 arêtes) associant de larges avenues et une structure de grille
  modifiée.
\item
  \textbf{Los Angeles} : Réseau de très grande taille (176 905 nœuds,
  619 100 arêtes) structuré sous forme de grille orthogonale régulière
  (Hgrid) avec de nombreuses voies d'accès rapides.
\item
  \textbf{Hanoï} : Réseau géant (473 209 nœuds, 1 656 200 arêtes)
  modélisant une mégapole asiatique caractérisée par une topologie
  irrégulière, des voies de service étroites et de nombreuses impasses.
\end{itemize}

\subsubsection{Protocole de simulation et variations paramétriques de la
demande}\label{protocole-de-simulation-et-variations-paramuxe9triques-de-la-demande}

Le corpus d'apprentissage est constitué à partir de \textbf{50
simulations d'échelle métropolitaine} entièrement exécutées dans notre
framework. Pour chaque ville, nous avons défini un protocole de
variation systématique des paramètres d'entrée afin de balayer tout le
spectre des régimes de trafic :

\begin{enumerate}
\def\labelenumi{\arabic{enumi}.}
\tightlist
\item
  \textbf{Volume cinématique} : Le nombre total de véhicules injectés
  dans le réseau varie de manière progressive (par exemple, de 1 000 à
  plus de 128 000 véhicules simultanés sur Paris) pour forcer le modèle
  à traverser les transitions de phase, depuis l'écoulement libre
  jusqu'à la congestion totale (\emph{gridlock}).
\item
  \textbf{Composition catégorielle de la flotte} : Pour chaque run, nous
  modifions la répartition des types de véhicules (voitures de tourisme,
  deux-roues motorisés, camions lourds et bus) afin de capturer
  l'impédance cinématique et le profil d'émission propre à chaque
  classe.
\item
  \textbf{Taux d'électrification découplés} : Nous faisons varier de
  manière indépendante le taux d'adoption des motorisations électriques
  (\(0\ %
  \) à \(15\ %
  \)) pour les voitures (\texttt{pct\_car\_ev}), les bus
  (\texttt{pct\_bus\_ev}), les camions (\texttt{pct\_truck\_ev}) et les
  motos (\texttt{pct\_moto\_ev}), agissant comme des modérateurs
  d'émissions directs.
\end{enumerate}

L'ensemble de ces simulations est exécuté en utilisant le mode de
simulation mésoscopique et le module de collecte rapide (\emph{Fast Data
Collection}) pour optimiser le temps de calcul et éviter les goulots
d'étranglement matériels (SWAP), accumulant ainsi une vérité terrain
simulée extrêmement robuste.

\subsubsection{Structuration des bases de données et modélisation des
régimes
météorologiques}\label{structuration-des-bases-de-donnuxe9es-et-moduxe9lisation-des-ruxe9gimes-muxe9tuxe9orologiques}

Plutôt que d'intégrer les conditions météorologiques (pluie, neige)
comme une simple variable d'entrée supplémentaire au sein d'un unique
modèle d'intelligence artificielle, nous optons pour une
\textbf{modélisation physique découplée}. En effet, la météo ne modifie
pas les rejets de \(CO_2\) de manière additive ou linéaire. Elle altère
profondément les équations cinématiques fondamentales au niveau
microscopique (adhérence des pneus, distances de sécurité, temps de
réaction, inattention), ce qui induit des perturbations macroscopiques
extrêmement non-linéaires qui dépendent de la géométrie de chaque ville
et du niveau d'adaptation régional des usagers.

Pour capturer cette complexité sans biaiser le métamodèle, nous
structurons notre apprentissage autour de \textbf{trois bases de données
d'apprentissage distinctes} générées par notre pipeline de simulation :

\begin{enumerate}
\def\labelenumi{\arabic{enumi}.}
\tightlist
\item
  \textbf{Le corpus ``Sol Sec'' (Clear/Optimal Weather)} : Il compile
  les simulations de référence exécutées sous des conditions optimales
  d'adhérence (\(\mu = 0.8\)) et d'attention standard des conducteurs
  (\(\sigma = 0.5\)).
\item
  \textbf{Le corpus ``Pluie'' (Rain Weather)} : Il compile les mêmes
  scénarios de simulations mais avec des paramètres physiques dégradés :
  coefficient d'adhérence réduit (\(\mu = 0.5\)) et imprécision
  comportementale accrue (\(\sigma\)) pour simuler la baisse de
  visibilité.
\item
  \textbf{Le corpus ``Neige'' (Snow Weather)} : Il compile les
  simulations exécutées sous conditions hivernales extrêmes :
  coefficient d'adhérence critique (\(\mu = 0.2\)) et forte inattention
  (\(\sigma\)).
\end{enumerate}

\paragraph{Impact physique de la météo sur la cinématique de
poursuite}\label{impact-physique-de-la-muxe9tuxe9o-sur-la-cinuxe9matique-de-poursuite}

Dans le micro-simulateur SUMO, ces perturbations météo agissent
directement sur le modèle comportemental de poursuite de véhicule
(Krauß). La distance de freinage sécuritaire est régie par la relation :
\[v_{safe}(t) = v_L(t) + \\frac{g(t) - v_L(t)\\tau}{\\frac{v(t) + v_L(t)}{2d} + \\tau}\]
Sous les intempéries, la décélération maximale \(d\) (fortement
dépendante de l'adhérence) subit une réduction structurelle de
\(30\ \%\) (sous la pluie) à \(70\ \%\) (sur la neige). Parallèlement,
l'imprécision de conduite \(\sigma\) (inattention) élargit l'intervalle
d'erreur stochastique \(\eta\), augmentant de fait le temps de réaction
\(\\tau\).

Ces dégradations cinématiques se traduisent macroscopiquement par un
affaissement de la vitesse moyenne globale (\(avg\_speed\_mps\)) et par
l'apparition d'ondes de choc d'arrêt-démarrage (cycles de
freinage-accélération). Ce sont ces cycles transitoires répétés qui
causent des surémissions de \(CO_2\) massives. Comme cet impact dépend
de l'accumulation géométrique (les zones de stockage et intersections
goulots d'étranglement amplifiant l'onde de choc), une simple constante
``météo'' dans le vecteur d'apprentissage de l'IA serait incapable de
prédire ces surcharges localisées. Il est donc indispensable d'exécuter
des simulations physiques complètes pour chaque régime météo afin de
générer la base de données propre à chaque climat.

\paragraph{Sociologie et disparités géographiques de
l'adaptation}\label{sociologie-et-disparituxe9s-guxe9ographiques-de-ladaptation}

L'analyse de sensibilité menée sur nos simulations met en lumière que la
réponse cinématique et environnementale face à une intempérie dépend de
l'habituation sociologique des usagers et de l'équipement des services
municipaux de la région ciblée :

\begin{itemize}
\tightlist
\item
  \textbf{Montagnes et Pré-alpes (Chamonix, Genève)} : Présentent une
  très forte adaptation. L'obligation d'équipements hivernaux (pneus
  neige) et l'efficacité des services de déneigement limitent la
  réduction de vitesse à seulement \(-5\ \%\\text{ à }-10\ \%\) sous la
  pluie, et \(-10\ \%\\text{ à }-15\ \%\) sous la neige.
\item
  \textbf{Plaines Européennes (Paris, Londres, Versailles)} : Présentent
  une habituation faible. L'absence d'équipements obligatoires et la
  saturation rapide des services de salage provoquent une baisse de
  vitesse de \(-10\ \%\\text{ à }-15\ \%\) sous la pluie, et un
  effondrement de \(-20\ \%\\text{ à }-35\ \%\) sous la neige (effet de
  panique urbaine).
\item
  \textbf{Climats Chauds \& Mousson (Hanoï, Casablanca, Le Caire)} :
  Caractérisés par une vulnérabilité extrême. Les pluies torrentielles
  de type mousson provoquent des inondations de chaussée entraînant des
  baisses de vitesse de \(-20\ \%\\text{ à }-30\ \%\). De plus, la neige
  y étant inédite, son apparition hypothétique paralyserait totalement
  le réseau (chute de vitesse supérieure à \(-50\ \%\) ou situation de
  blocage complet \emph{gridlock}).
\item
  \textbf{Grilles Nord-Américaines (Los Angeles)} : Les précipitations y
  sont très rares. En cas d'intempérie, le manque d'expérience des
  usagers provoque une multiplication des accidents mineurs et des
  freinages préventifs abrupts sur les autoroutes urbaines, causant une
  baisse de vitesse disproportionnée de \(-15\ \%\\text{ à }-20\ \%\)
  sous la pluie (la neige n'y étant pas applicable).
\item
  \textbf{Maillage Océanien (Sydney, Auckland)} : Présentent une
  sensibilité standard, avec une réduction d'environ \(-10\ \%\) sous la
  pluie due aux contraintes classiques de visibilité et de prudence.
\end{itemize}

En segmentant ainsi notre apprentissage en trois modèles XGBoost dédiés
(Sec, Pluie, Neige) entraînés chacun sur leur base de données de
simulations respectives (contenant 50 simulations d'échelle
métropolitaine par base), nous parvenons à capturer de manière endogène
l'impact non-linéaire et régionalisé des intempéries sur la pollution.

\newpage

\section{DEUXIÈME PARTIE : FONDATIONS TECHNIQUES ET THÉORIQUES DES
SYSTÈMES DE SIMULATION ET DE LA PRÉDICTION
TOPOLOGIQUE}\label{deuxiuxe8me-partie-fondations-techniques-et-thuxe9oriques-des-systuxe8mes-de-simulation-et-de-la-pruxe9diction-topologique}

\subsection{Chapitre 5 : Le moteur de micro-simulation physique
SUMO}\label{chapitre-5-le-moteur-de-micro-simulation-physique-sumo}

\subsubsection{Abstraction topologique de la voirie (Nœuds, arêtes,
connecteurs
géométriques)}\label{abstraction-topologique-de-la-voirie-nux153uds-aruxeates-connecteurs-guxe9omuxe9triques}

Le progiciel SUMO modélise les réseaux de transport sous forme de
réseaux logiques basés sur la théorie des graphes orientés. Dans ce
formalisme, chaque intersection physique est représentée par un nœud
unique doté d'une géométrie polygonale décrivant sa surface de jonction.
Les tronçons routiers reliant les nœuds sont modélisés par des arêtes,
subdivisées en une ou plusieurs voies de circulation (\emph{lanes}).

Chaque voie possède des attributs géométriques et comportementaux
stricts : * Une polyligne tridimensionnelle décrivant son axe central. *
Une largeur constante (généralement fixée à 3,2 mètres pour les voies
urbaines standards). * Une liste de classes de véhicules autorisées (ex:
\texttt{passenger}, \texttt{taxi}, \texttt{motorcycle}, \texttt{bus},
\texttt{delivery}). * Une vitesse limite supérieure déterminant la
vitesse de référence des agents.

La transition entre deux arêtes consécutives s'effectue via des
\textbf{connecteurs géométriques} (\emph{connections}) définis à
l'intérieur des nœuds. Ces connecteurs lient précisément une voie de
l'arête d'approche à une voie de l'arête de sortie. Ils supportent les
règles de priorité (ex: céder le passage, priorité absolue) et les
configurations de signalisation dynamique (phases de feux).

\subsubsection{Modèle de poursuite cinématique de Krauß (Vitesse
sécuritaire, temps de réaction, bruit
stochastique)}\label{moduxe8le-de-poursuite-cinuxe9matique-de-krauuxdf-vitesse-suxe9curitaire-temps-de-ruxe9action-bruit-stochastique}

Pour reproduire le mouvement individuel des véhicules le long des
arêtes, SUMO implémente par défaut le modèle comportemental de poursuite
de véhicule développé par \textbf{Krauß}. Ce modèle cinématique calcule
à chaque pas de temps la vitesse optimale d'un véhicule suiveur pour
éviter toute collision avec le véhicule leader, même si ce dernier
décélère brutalement.

Soit un véhicule suiveur \(F\) caractérisé par sa position \(x(t)\) et
sa vitesse \(v(t)\), circulant derrière un véhicule leader \(L\)
caractérisé par sa position \(x_L(t)\) et sa vitesse \(v_L(t)\).
L'intervalle spatial libre (ou gap) séparant les deux véhicules est
défini par : \[g(t) = x_L(t) - x(t) - l_L\] où \(l_L\) est la longueur
physique du véhicule leader.

Le modèle détermine la vitesse sécuritaire \(v_{safe}(t)\) par la
relation suivante :
\[v_{safe}(t) = v_L(t) + \frac{g(t) - v_L(t)\tau}{\frac{v(t) + v_L(t)}{2d} + \tau}\]
où \(\tau\) représente le temps de réaction du conducteur, et \(d\) sa
capacité de décélération maximale.

La vitesse théorique souhaitée pour le pas de temps suivant,
\(v_{target}(t)\), est le minimum parmi les contraintes physiques et
légales :
\[v_{target}(t) = \min\left( V_{max}, v(t) + a \cdot \Delta t, v_{safe}(t) \right)\]
où \(V_{max}\) est la vitesse limite de la voie, et \(a\) est la
capacité d'accélération maximale du véhicule.

Enfin, pour introduire la variabilité des comportements humains (retards
de réaction légers, imprécisions de contrôle), une perturbation
stochastique négative \(\eta\) est soustraite de la vitesse cible pour
obtenir la vitesse finale appliquée à l'agent :
\[v(t + \Delta t) = \max\left( 0, v_{target}(t) - \eta \right)\] où
\(\eta\) est une variable aléatoire distribuée uniformément sur
l'intervalle \([0, \sigma \cdot a \cdot \Delta t]\), le paramètre
\(\sigma \in [0, 1]\) caractérisant le degré d'inattention du
conducteur.

\subsubsection{Pipeline de traitement topologique : Compilation des
réseaux OSM, nettoyage par
netconvert}\label{pipeline-de-traitement-topologique-compilation-des-ruxe9seaux-osm-nettoyage-par-netconvert}

La conversion de données topologiques brutes issues d'OpenStreetMap
(OSM) en réseaux de simulation exploitables par SUMO nécessite
l'exécution d'un pipeline de traitement rigoureux via l'outil
\texttt{netconvert}.

Ce pipeline effectue les opérations de nettoyage suivantes : 1.
\textbf{Uniformisation géométrique :} Re-projection des coordonnées
géographiques sphériques (WGS84) vers un système de coordonnées
cartésiennes planes UTM (Universal Transverse Mercator), permettant
d'effectuer les calculs de distance et de vitesse en mètres et mètres
par seconde. 2. \textbf{Simplification des nœuds :} Fusion des grappes
d'intersections complexes et des micro-carrefours d'OSM en nœuds
logiques uniques. Cette opération élimine les segments de voirie
artificiellement courts (inférieurs à 5 mètres) qui s'avèrent
impossibles à gérer pour les modèles cinématiques de car-following. 3.
\textbf{Nettoyage des connexions redondantes :} Suppression des voies
d'accès et des bretelles d'insertion non connectées ou mal définies dans
la base de données source, évitant la création de trajectoires
conflictuelles anormales.

\subsubsection{Connexité réseau : Algorithme de Tarjan pour l'extraction
des Composantes Fortement Connexes
(SCC)}\label{connexituxe9-ruxe9seau-algorithme-de-tarjan-pour-lextraction-des-composantes-fortement-connexes-scc}

L'intégrité topologique d'un réseau routier de simulation dépend de sa
connexité globale. En raison de la présence de règles de circulation
directionnelles (sens uniques) et de potentielles erreurs de
numérisation dans les données géospatiales d'entrée, un graphe routier
orienté brut comporte souvent des composantes déconnectées du flux
principal.

Pour écarter ce problème, notre pipeline met en œuvre
l'\textbf{algorithme des Composantes Fortement Connexes (SCC)} de
Tarjan.

Soit un graphe orienté \(G = (V, E)\). Une composante fortement connexe
de \(G\) est un sous-graphe maximal \(G' = (V', E')\) dans lequel il
existe un chemin orienté reliant chaque paire de nœuds de manière
bidirectionnelle :
\[\forall (u, v) \in V'^2, \quad u \rightsquigarrow v \quad \text{et} \quad v \rightsquigarrow u\]

L'algorithme de Tarjan utilise un parcours en profondeur (DFS) pour
identifier les cycles et partitionner le graphe en composantes fortement
connexes en une seule passe, avec une complexité temporelle optimale :
\[\mathcal{O}(|V| + |E|)\]

Une fois la partition \(\mathcal{C} = \{G_1, G_2, \dots, G_k\}\)
obtenue, le pipeline extrait la composante principale possédant le plus
grand cardinal de nœuds :
\[G_{SCC_{max}} = \arg\max_{G_i \in \mathcal{C}} |V_i|\]

Tous les nœuds et arêtes n'appartenant pas à cette composante principale
sont supprimés du réseau de simulation final. Ce traitement garantit que
chaque véhicule injecté sur le réseau dispose d'un chemin physique
valide pour atteindre n'importe quel point de destination, éliminant les
blocages d'agents et les échecs de routage dynamique.

\subsubsection{Routage dynamique : Filtre de distance minimale pour
l'elimination des micro-trajets
parasitaires}\label{routage-dynamique-filtre-de-distance-minimale-pour-lelimination-des-micro-trajets-parasitaires}

Lors de la phase de génération automatique de la demande de trafic
(synthèse des trajets), les points d'origine et de destination sont
distribués aléatoirement sur le graphe épuré. Pour éviter l'apparition
de micro-trajets (déplacements de moins de 300 mètres reliant des
intersections adjacentes), nous implémentons une contrainte de distance
minimale lors de la phase de routage.

Soit un couple origine-destination \((o, d) \in V^2\) sélectionné pour
générer le trajet d'un agent. Le trajet n'est validé et écrit dans le
fichier final de demande que s'il satisfait la condition suivante :
\[D_{Euclidienne}(o, d) = \sqrt{(x_d - x_o)^2 + (y_d - y_o)^2} \ge 300 \text{ mètres}\]

Cette contrainte force le planificateur d'itinéraires à rejeter les
trajets de très courte distance. En éliminant ces mouvements
parasitaires qui se limitent à des phases d'insertion-sortie immédiates,
le filtre garantit que l'ensemble de la flotte simulée s'insère dans les
flux de transit principaux du réseau.

\newpage

\subsection{Chapitre 6 : Théorie mathématique du trafic et de la
topologie
spectrale}\label{chapitre-6-thuxe9orie-mathuxe9matique-du-trafic-et-de-la-topologie-spectrale}

L'évaluation de la résilience et des émissions de carbone au sein d'un
réseau viaire repose traditionnellement sur des modèles dynamiques
microscopiques ou sur des approximations macroscopiques fluides. Bien
que ces méthodes offrent une précision locale indéniable, elles
souffrent d'une lenteur computationnelle rébarbative et n'offrent que
peu d'indications théoriques a priori sur la vulnérabilité intrinsèque
de la structure. Dans ce chapitre, nous développons un formalisme
mathématique rigoureux qui extrait la ``signature urbaine'' d'une ville
à travers l'analyse spectrale et non-normale de sa matrice d'adjacence.
Ce cadre théorique permet de quantifier l'instabilité potentielle du
trafic et d'identifier les fragilités topologiques structurelles avant
même toute simulation cinématique.

\subsubsection{Formalisation matricielle du réseau
viaire}\label{formalisation-matricielle-du-ruxe9seau-viaire}

Pour modéliser mathématiquement le réseau routier, nous le représentons
sous la forme d'un graphe orienté et pondéré \(G = (V, E)\), où : *
\(V = \{v_1, v_2, \dots, v_n\}\) désigne l'ensemble des nœuds
(\(n = |V|\)), représentant les intersections physiques du réseau. *
\(E \subset V \times V\) désigne l'ensemble des arêtes orientées
(\(m = |E|\)), modélisant les tronçons routiers à sens unique reliant
ces intersections.

La connectivité et l'impédance physique du réseau sont codées dans la
\textbf{matrice d'adjacence pondérée} \(A \in \mathbb{R}^{n \times n}\).
Contrairement aux graphes non orientés simples où la matrice d'adjacence
est symétrique et binaire, la représentation réaliste d'un réseau urbain
impose une double complexité : 1. \textbf{Asymétrie structurelle :}
L'existence de sens uniques et de priorités de passage implique que
l'existence d'un arc \(v_i \to v_j\) n'entraîne pas celle de l'arc
\(v_j \to v_i\). Ainsi, \(A_{ij} \neq A_{ji}\). 2. \textbf{Pondération
physique :} Chaque coefficient \(A_{ij}\) n'indique pas une simple
connexion, mais est défini comme une fonction de la capacité physique du
tronçon routier. Pour un tronçon reliant le nœud \(i\) au nœud \(j\), la
pondération est donnée par :
\[A_{ij} = \begin{cases} \frac{L_{ij} \cdot C_{ij}}{W_{ij}} & \text{si } (v_i, v_j) \in E \\ 0 & \text{sinon} \end{cases}\]
où \(L_{ij}\) représente la longueur du tronçon (en mètres), \(C_{ij}\)
le nombre de voies de circulation, et \(W_{ij}\) la vitesse maximale
autorisée (en m/s). Le rapport \(\frac{L_{ij}}{W_{ij}}\) représente le
\textbf{temps de parcours libre} de l'arête. En le multipliant par
\(C_{ij}\), on intègre sa capacité d'accumulation. Un boulevard large
aura ainsi un poids spectral supérieur à une rue étroite.

Pour illustrer ce formalisme, considérons un exemple minimal de réseau
routier à 4 nœuds représentant une intersection en boucle fermée avec un
goulot d'étranglement :

\begin{verbatim}
       (v1) ---------> (v2)
        ^               |
        |               | (Goulot d'étranglement)
        |               v
       (v4) <--------- (v3)
\end{verbatim}

Dans ce modèle simplifié, la matrice d'adjacence orientée non pondérée
\(A_{ex}\) s'écrit : \[A_{ex} = \begin{pmatrix}
0 & 1 & 0 & 0 \\
0 & 0 & 1 & 0 \\
0 & 0 & 0 & 0 \\
1 & 0 & 1 & 0
\end{pmatrix}\] L'absence de symétrie de cette matrice est évidente
(\(A_{ex} \neq A_{ex}^T\)).

\subsubsection{Le concept de non-normalité et amplification
transitoire}\label{le-concept-de-non-normalituxe9-et-amplification-transitoire}

Une matrice carrée \(A\) est dite \textbf{normale} si et seulement s'il
commute avec sa transposée, soit \(A A^T = A^T A\). Dans le cas des
réseaux routiers réels orientés, cette relation n'est jamais vérifiée :
la matrice d'adjacence \(A\) est intrinsèquement \textbf{non-normale}
(\(A A^T \neq A^T A\)).

La non-symétrie de la matrice implique sa non-normalité. Pour quantifier
ce phénomène, nous introduisons \textbf{l'indice d'asymétrie}
\(\alpha(G)\) : \[\alpha(G) = 1.0 - \frac{|E_{bidirectionnel}|}{|E|}\]
Où \(|E_{bidirectionnel}|\) désigne le nombre d'arêtes admettant un arc
de retour identique. Un indice \(\alpha(G) \approx 1\) décrit un réseau
purement asymétrique (sens uniques dominants), tandis qu'un indice
proche de 0 décrit un réseau symétrique (double sens de circulation sur
tous les axes).

La mesure de cette non-normalité est quantifiée analytiquement par la
\textbf{norme de Frobenius du commutateur} :
\[\Delta(A) = \| A A^T - A^T A \|_F = \sqrt{\text{Tr}\left( (A A^T - A^T A)^H (A A^T - A^T A) \right)}\]

Cette non-normalité a des implications physiques profondes sur la
dynamique du trafic. Pour une matrice normale, les vecteurs propres
forment une base orthonormale, et toute perturbation décroît de manière
monotone si le système est stable. Dans un système non-normal, les
vecteurs propres ne sont plus orthogonaux et peuvent devenir presque
colinéaires, provoquant un \textbf{phénomène d'amplification
transitoire} (l'effet ``coup de bélier''). Une infime perturbation à une
intersection (ex. blocage d'un carrefour pendant 15 secondes) excite ces
modes non-orthogonaux (voir Figures D.5 et D.6, Annexe D). Avant de se
dissiper, la perturbation s'amplifie de manière géométrique à court
terme, provoquant des congestions diffuses inattendues et des
surémissions massives de \(CO_2\) causées par les cycles d'arrêt et de
réaccélération.

\subsubsection{\texorpdfstring{Le Rayon Spectral (\(\rho\)) et le
Théorème de
Perron-Frobenius}{Le Rayon Spectral (\textbackslash rho) et le Théorème de Perron-Frobenius}}\label{le-rayon-spectral-rho-et-le-thuxe9oruxe8me-de-perron-frobenius}

Le spectre d'une matrice, noté \(\sigma(A)\), regroupe ses valeurs
propres complexes \(\lambda_i \in \mathbb{C}\) résolvant
\(\det(\lambda I - A) = 0\). Le \textbf{rayon spectral} \(\rho(A)\)
correspond à la borne supérieure du module des valeurs propres :
\[\rho(A) = \max_{\lambda \in \sigma(A)} |\lambda|\]

Puisque les coefficients \(A_{ij}\) de notre matrice d'adjacence
pondérée sont strictement non-négatifs (\(A_{ij} \ge 0\)), nous pouvons
appliquer le \textbf{théorème de Perron-Frobenius} si le graphe est
fortement connexe (garanti par l'extraction des SCC) : 1. Le rayon
spectral \(\rho(A)\) est lui-même une valeur propre de \(A\), simple et
strictement positive (\(\rho(A) > 0\)). 2. Il existe un vecteur propre à
droite \(v_{PF}\) associé à \(\rho(A)\) dont toutes les composantes sont
strictement positives (\(v_{PF} > 0\)), appelé vecteur de
Perron-Frobenius. 3. Cette valeur propre domine toutes les autres :
\(\forall \lambda \in \sigma(A) \setminus \{\rho(A)\}, \ |\lambda| \le \rho(A)\).

Physiquement, le rayon spectral \(\rho(A)\) quantifie la
\textbf{capacité globale de diffusion} ou de transit du réseau. Une
valeur propre \(\rho\) élevée traduit un réseau doté d'axes structurants
dominants (corridors). Le vecteur \(v_{PF}\) identifie de manière
déterministe les zones où convergent naturellement les flux.

\subsubsection{\texorpdfstring{La Constante de Kreiss (\(K\)) et la
dynamique de
crise}{La Constante de Kreiss (K) et la dynamique de crise}}\label{la-constante-de-kreiss-k-et-la-dynamique-de-crise}

Pour quantifier rigoureusement la sensibilité d'un réseau non-normal aux
amplifications transitoires et modéliser son instabilité dynamique, nous
introduisons la \textbf{constante de Kreiss} \(K(A)\). Soit \(A\) une
matrice stable (\(\rho(A) < 1\)). La constante de Kreiss est définie par
: \[K(A) = \sup_{|z| > 1} (|z| - 1) \left\| (zI - A)^{-1} \right\|_2\]
où \(\|\cdot\|_2\) désigne la norme matricielle induite (norme
spectrale).

Le théorème des matrices de Kreiss établit des bornes strictes reliant
cette constante à l'amplification transitoire maximale de la puissance
de la matrice :
\[K(A) \le \sup_{k \ge 0} \left\| A^k \right\|_2 \le e \cdot n \cdot K(A)\]
où \(n\) is la dimension de la matrice.

En ingénierie du trafic, la constante de Kreiss agit comme un
\textbf{détecteur de nervosité} ou de fragilité structurelle du réseau.
Une constante élevée indique qu'une petite perturbation (un accident, un
feu trop long) va provoquer un effet domino foudroyant. C'est une
variable clé pour la prédiction de la pollution car une ville à
constante de Kreiss élevée aura une pollution instable et des pics très
forts en heure de pointe.

\subsubsection{\texorpdfstring{Les Normes de Hardy \(H_2\) et
\(H_\infty\)}{Les Normes de Hardy H\_2 et H\_\textbackslash infty}}\label{les-normes-de-hardy-h_2-et-h_infty}

En modélisant le réseau routier comme un filtre dynamique linéaire
entrée-sortie (où l'entrée est le flux d'injection des véhicules et la
sortie la congestion), nous évaluons sa robustesse via les normes
\(H_2\) et \(H_\infty\) de sa fonction de transfert
\(T(z) = (zI - A)^{-1}\) : 1. \textbf{La Norme \(H_\infty\) (Pire
scénario d'amplification)} :
\[\|T\|_{H_\infty} = \sup_{|z| > 1} \left\| (zI - A)^{-1} \right\|_2 = \sup_{\theta \in [0, 2\pi]} \sigma_{max}\left( (e^{i\theta}I - A)^{-1} \right)\]
Elle caractérise le niveau de pollution ``de base'' inévitable dû à la
géométrie de la ville (gain maximal stabilisé). 2. \textbf{La Norme
\(H_2\) (Énergie de perturbation stockée)} :
\[\|T\|_{H_2} = \left( \sum_{k=0}^{\infty} \left\| A^k \right\|_F^2 \right)^{1/2}\]
Elle caractérise la \textbf{mémoire temporelle de la congestion} : le
temps que met la pollution à s'évacuer après un pic d'affluence.

\subsubsection{Théorie des perturbations de Kato et Loi de
Contrôle}\label{thuxe9orie-des-perturbations-de-kato-et-loi-de-contruxf4le}

Dans le cadre de l'optimisation des réseaux urbains, une question
centrale se pose : comment modifier la structure du graphe pour
minimiser l'apparition des congestions et la pollution associée sans
avoir à recalculer intégralement le spectre de la matrice d'adjacence
(ce qui est extrêmement coûteux pour des réseaux de taille
métropolitaine) ?

Pour répondre à cela, nous modélisons les modifications d'infrastructure
comme des perturbations de la matrice d'adjacence pondérée \(A\) sous la
forme : \[\delta A = \epsilon B\] où : * \(\epsilon \in \mathbb{R}^+\)
est un paramètre d'échelle infinitésimal régissant l'intensité globale
de la modification. * \(B \in \mathbb{R}^{n \times n}\) désigne la
\textbf{matrice de perturbation (ou matrice de contrôle topologique)}.
Chaque élément \(B_{ij}\) quantifie l'action d'aménagement local sur le
tronçon orienté reliant le nœud \(i\) au nœud \(j\) : - \(B_{ij} > 0\)
correspond à une expansion de capacité physique (e.g., ajout d'une voie
de circulation, recalibrage géométrique d'une intersection) ou à des
aménagements fluidifiant l'écoulement (e.g., instauration d'une onde
verte de feux tricolores, hausse de la vitesse limite). - \(B_{ij} < 0\)
correspond à une réduction de capacité (e.g., piétonnisation d'un axe,
réduction de la vitesse réglementaire, aménagement de pistes cyclables
ou de couloirs de bus exclusifs). - \(B_{ij} = 0\) pour tous les liens
dont la géométrie reste inchangée ou pour lesquels aucun aménagement
n'est physiquement réalisable.

En s'appuyant sur la \textbf{théorie des perturbations de Kato} (Kato,
1995), nous pouvons évaluer l'impact analytique d'une telle perturbation
sur les valeurs propres \(\lambda_i\) du système. Pour une valeur propre
simple de \(A\), ses dérivées de premier et second ordre par rapport à
la perturbation sont données par :

\begin{enumerate}
\def\labelenumi{\arabic{enumi}.}
\tightlist
\item
  \textbf{Dérivée Première (Sensibilité linéaire)} :
  \[\lambda_i^{(1)} = \frac{w_i^T B v_i}{w_i^T v_i}\] où \(v_i\) et
  \(w_i\) désignent respectivement les vecteurs propres à droite et à
  gauche associés à \(\lambda_i\).
\item
  \textbf{Dérivée Seconde (Couplage non-linéaire du second ordre)} :
  \[\lambda_i^{(2)} = w_i^T B S_i B v_i\] Où \(S_i\) représente la
  résolvante réduite définie par
  \(S_i = \lim_{z \to \lambda_i} (zI - A)^{-1}(I - P_i)\), avec
  \(P_i = \frac{v_i w_i^T}{w_i^T v_i}\) le projecteur spectral associé.
  Cette résolvante modélise la propagation indirecte de la perturbation
  à travers les chemins alternatifs du graphe.
\end{enumerate}

\textbf{Loi de contrôle topologique :} L'objectif d'un planificateur
urbain est de concevoir une stratégie d'aménagement optimale \(B^*\)
appartenant à l'espace des contrôles admissibles \(\mathcal{B}\) pour
minimiser le rayon spectral \(\rho(A + \epsilon B)\) (et donc réduire
l'impédance de transit globale et le risque de congestion sous charge).
En exploitant le théorème de Perron-Frobenius, cette loi de contrôle
s'exprime par le problème de minimisation sous contraintes suivant :
\[B^* = \arg\min_{B \in \mathcal{B}} \rho(A + \epsilon B) \approx \arg\min_{B \in \mathcal{B}} \left( \lambda_{PF} + \epsilon \frac{w_{PF}^T B v_{PF}}{w_{PF}^T v_{PF}} + \epsilon^2 w_{PF}^T B S_{PF} B v_{PF} \right)\]
où \(\lambda_{PF}\), \(v_{PF}\) et \(w_{PF}\) désignent les valeurs et
vecteurs propres dominants de Perron-Frobenius associés au graphe
initial stable.

L'espace des contrôles admissibles \(\mathcal{B}\) est structuré par des
contraintes opérationnelles réelles : 1. \textbf{Contrainte budgétaire
et d'échelle (Sparsité et volume d'action)} :
\[\sum_{(i,j) \in E} |B_{ij}| \le C_{budget}\] Le coût total des travaux
de restructuration ou d'ajout de capacité est borné supérieurement. 2.
\textbf{Limites géométriques locales} :
\[B_{ij} \le B_{max} \quad \forall (i,j) \in E\] La largeur de voirie et
le foncier urbain limitent physiquement le nombre maximal de voies qu'il
est possible d'ajouter sur un tronçon donné. 3. \textbf{Contraintes
patrimoniales et géographiques (Sparsité stricte)} :
\[B_{ij} = 0 \quad \forall (i,j) \notin E_{modifiable}\] Certaines zones
(hyper-centres historiques protégés, ponts critiques) ne peuvent subir
de modifications géométriques, imposant des coefficients nuls dans la
matrice de contrôle.

\textbf{Interprétation physique de la loi de contrôle :} Le produit
matriciel du premier ordre
\(w_{PF}^T B v_{PF} = \sum_{i,j} w_{PF, i} B_{ij} v_{PF, j}\) montre que
l'efficacité d'un aménagement sur l'arête \((i, j)\) dépend du couplage
entre la centralité de diffusion du nœud source (mesurée par la
composante à gauche \(w_{PF, i}\)) et l'attractivité du nœud cible
(mesurée par la composante à droite \(v_{PF, j}\)). Ainsi, pour
maximiser la résilience du trafic, les investissements doivent se
concentrer en priorité sur les liaisons reliant des intersections
hautement distributrices à des intersections réceptrices critiques.

Le terme de second ordre intègre quant à lui les transferts de
congestion : il empêche le déplacement simple du goulot d'étranglement
vers des tronçons adjacents en pénalisant les perturbations qui
surchargent la résolvante réduite \(S_{PF}\).

\subsubsection{Données Expérimentales d'Analyse Topologique et
Spectrale}\label{donnuxe9es-expuxe9rimentales-danalyse-topologique-et-spectrale}

L'analyse systématique des réseaux routiers de notre base de données
(présente dans \texttt{IA\_report/MATHEMATICAL\_TOPOLOGY\_REPORT.md}) a
permis d'extraire les caractéristiques topologiques et spectrales
réelles présentées dans les tableaux ci-dessous.

\paragraph{Caractéristiques Topologiques Générales des
Villes}\label{caractuxe9ristiques-topologiques-guxe9nuxe9rales-des-villes}

{\def\LTcaptype{none} % do not increment counter
\begin{longtable}[]{@{}
  >{\raggedright\arraybackslash}p{(\linewidth - 16\tabcolsep) * \real{0.1111}}
  >{\raggedright\arraybackslash}p{(\linewidth - 16\tabcolsep) * \real{0.1111}}
  >{\raggedleft\arraybackslash}p{(\linewidth - 16\tabcolsep) * \real{0.1111}}
  >{\raggedleft\arraybackslash}p{(\linewidth - 16\tabcolsep) * \real{0.1111}}
  >{\raggedleft\arraybackslash}p{(\linewidth - 16\tabcolsep) * \real{0.1111}}
  >{\raggedleft\arraybackslash}p{(\linewidth - 16\tabcolsep) * \real{0.1111}}
  >{\raggedleft\arraybackslash}p{(\linewidth - 16\tabcolsep) * \real{0.1111}}
  >{\raggedleft\arraybackslash}p{(\linewidth - 16\tabcolsep) * \real{0.1111}}
  >{\raggedleft\arraybackslash}p{(\linewidth - 16\tabcolsep) * \real{0.1111}}@{}}
\toprule\noalign{}
\begin{minipage}[b]{\linewidth}\raggedright
city
\end{minipage} & \begin{minipage}[b]{\linewidth}\raggedright
origin
\end{minipage} & \begin{minipage}[b]{\linewidth}\raggedleft
nodes
\end{minipage} & \begin{minipage}[b]{\linewidth}\raggedleft
edges
\end{minipage} & \begin{minipage}[b]{\linewidth}\raggedleft
density
\end{minipage} & \begin{minipage}[b]{\linewidth}\raggedleft
avg\_degree
\end{minipage} & \begin{minipage}[b]{\linewidth}\raggedleft
asymmetry\_index
\end{minipage} & \begin{minipage}[b]{\linewidth}\raggedleft
sources\_ratio
\end{minipage} & \begin{minipage}[b]{\linewidth}\raggedleft
sinks\_ratio
\end{minipage} \\
\midrule\noalign{}
\endhead
\bottomrule\noalign{}
\endlastfoot
\textbf{nairobi} & Africa & 40 685 & 89 950 & 5.43e-05 & 2.21 & -0.44 &
0.002 & 0.002 \\
\textbf{marseille} & Europe & 17 035 & 34 858 & 1.20e-04 & 2.05 & -0.14
& 0.021 & 0.020 \\
\textbf{cairo} & Africa & 13 095 & 29 575 & 1.72e-04 & 2.26 & -0.33 &
0.005 & 0.005 \\
\textbf{london} & Europe & 12 415 & 28 623 & 1.86e-04 & 2.31 & -0.32 &
0.021 & 0.014 \\
\textbf{casablanca} & Africa & 9 493 & 22 098 & 2.45e-04 & 2.33 & -0.28
& 0.015 & 0.015 \\
\textbf{berlin} & Europe & 8 855 & 21 043 & 2.68e-04 & 2.38 & -0.40 &
0.007 & 0.006 \\
\textbf{amsterdam} & Europe & 7 088 & 16 073 & 3.20e-04 & 2.27 & -0.13 &
0.035 & 0.029 \\
\textbf{lyon} & Europe & 6 356 & 11 661 & 2.89e-04 & 1.83 & 0.18 & 0.033
& 0.030 \\
\textbf{los\_angeles} & North\_America & 5 947 & 13 810 & 3.91e-04 &
2.32 & -0.35 & 0.008 & 0.009 \\
\textbf{madrid} & Europe & 5 651 & 10 567 & 3.31e-04 & 1.87 & 0.37 &
0.041 & 0.038 \\
\textbf{geneva} & Europe & 5 634 & 12 314 & 3.88e-04 & 2.19 & -0.23 &
0.021 & 0.019 \\
\textbf{paris} & Europe & 5 199 & 9 904 & 3.66e-04 & 1.90 & 0.19 & 0.048
& 0.042 \\
\textbf{sydney} & Oceania & 4 823 & 10 483 & 4.51e-04 & 2.17 & -0.25 &
0.024 & 0.021 \\
\textbf{dubai} & Asia\_Middle\_East & 3 841 & 6 476 & 4.39e-04 & 1.69 &
0.37 & 0.039 & 0.037 \\
\textbf{hanoi} & Asia\_Middle\_East & 3 195 & 7 256 & 7.11e-04 & 2.27 &
-0.28 & 0.015 & 0.013 \\
\textbf{strasbourg} & Europe & 2 945 & 6 070 & 7.00e-04 & 2.06 & -0.27 &
0.036 & 0.034 \\
\textbf{buenos\_aires} & South\_America & 2 820 & 5 637 & 7.09e-04 &
2.00 & 0.49 & 0.018 & 0.017 \\
\textbf{versailles} & Europe & 1 794 & 3 686 & 1.15e-03 & 2.05 & -0.08 &
0.046 & 0.035 \\
\textbf{rio\_de\_janeiro} & South\_America & 1 628 & 3 092 & 1.17e-03 &
1.90 & 0.16 & 0.015 & 0.014 \\
\textbf{chamonix} & Europe & 848 & 1 896 & 2.64e-03 & 2.24 & -0.35 &
0.009 & 0.012 \\
\textbf{monaco} & Europe & 672 & 1 286 & 2.85e-03 & 1.91 & 0.00 & 0.039
& 0.039 \\
\end{longtable}
}

\paragraph{Métriques Spectrales de Non-Normalité (Matrices
Non-Pondérées)}\label{muxe9triques-spectrales-de-non-normalituxe9-matrices-non-ponduxe9ruxe9es}

{\def\LTcaptype{none} % do not increment counter
\begin{longtable}[]{@{}
  >{\raggedright\arraybackslash}p{(\linewidth - 10\tabcolsep) * \real{0.1667}}
  >{\raggedleft\arraybackslash}p{(\linewidth - 10\tabcolsep) * \real{0.1667}}
  >{\raggedleft\arraybackslash}p{(\linewidth - 10\tabcolsep) * \real{0.1667}}
  >{\raggedleft\arraybackslash}p{(\linewidth - 10\tabcolsep) * \real{0.1667}}
  >{\raggedleft\arraybackslash}p{(\linewidth - 10\tabcolsep) * \real{0.1667}}
  >{\raggedleft\arraybackslash}p{(\linewidth - 10\tabcolsep) * \real{0.1667}}@{}}
\toprule\noalign{}
\begin{minipage}[b]{\linewidth}\raggedright
city
\end{minipage} & \begin{minipage}[b]{\linewidth}\raggedleft
non\_normalness
\end{minipage} & \begin{minipage}[b]{\linewidth}\raggedleft
spectral\_radius
\end{minipage} & \begin{minipage}[b]{\linewidth}\raggedleft
sigma\_max
\end{minipage} & \begin{minipage}[b]{\linewidth}\raggedleft
h2\_norm
\end{minipage} & \begin{minipage}[b]{\linewidth}\raggedleft
kreiss\_constant
\end{minipage} \\
\midrule\noalign{}
\endhead
\bottomrule\noalign{}
\endlastfoot
\textbf{nairobi} & 137.04 & 4.717 & 4.720 & 419.80 & 8.49 \\
\textbf{marseille} & 207.55 & 4.149 & 4.149 & 247.51 & 24.17 \\
\textbf{cairo} & 155.40 & 4.571 & 4.571 & 236.32 & 32.50 \\
\textbf{london} & 162.56 & 4.378 & 4.378 & 231.54 & 0.00 \\
\textbf{casablanca} & 149.26 & 5.545 & 5.545 & 202.30 & 17.89 \\
\textbf{berlin} & 107.52 & 4.272 & 4.272 & 201.32 & 8.49 \\
\textbf{amsterdam} & 170.33 & 4.553 & 4.553 & 167.25 & 30.46 \\
\textbf{lyon} & 140.53 & 4.110 & 4.251 & 134.16 & 21.54 \\
\textbf{los\_angeles} & 87.77 & 4.388 & 4.423 & 161.94 & 14.70 \\
\textbf{madrid} & 159.09 & 4.149 & 4.149 & 122.05 & 29.93 \\
\textbf{geneva} & 123.09 & 4.000 & 4.000 & 149.58 & 19.39 \\
\textbf{paris} & 143.29 & 3.939 & 4.047 & 123.38 & 8.49 \\
\textbf{sydney} & 104.20 & 4.202 & 4.202 & 138.37 & 20.40 \\
\textbf{dubai} & 109.66 & 4.123 & 4.279 & 96.33 & 30.33 \\
\textbf{hanoi} & 92.63 & 4.615 & 4.615 & 115.79 & 34.41 \\
\textbf{strasbourg} & 75.71 & 4.000 & 4.000 & 105.71 & 30.98 \\
\textbf{buenos\_aires} & 121.23 & 4.000 & 4.106 & 86.94 & 45.43 \\
\textbf{versailles} & 76.62 & 4.202 & 4.202 & 79.44 & 38.57 \\
\textbf{rio\_de\_janeiro} & 71.64 & 3.696 & 3.696 & 69.47 & 32.00 \\
\textbf{chamonix} & 36.19 & 3.696 & 3.901 & 60.00 & 33.59 \\
\textbf{monaco} & 45.59 & 3.604 & 3.604 & 46.26 & 62.55 \\
\end{longtable}
}

\newpage

\subsection{Chapitre 7 : Le moteur d'intelligence artificielle prédictif
direct}\label{chapitre-7-le-moteur-dintelligence-artificielle-pruxe9dictif-direct}

Pour s'affranchir définitivement du coût computationnel prohibitif des
simulateurs physiques multi-agents comme SUMO, ce chapitre présente la
conception, la formulation et la validation expérimentale de notre
métamodèle d'intelligence artificielle. En s'appuyant sur la ``signature
urbaine'' spectrale et topologique développée au Chapitre 6, ce moteur
réalise des prédictions instantanées des émissions de \(CO_2\)
directement depuis la structure du réseau et la charge de trafic sans
nécessiter d'ajustements intermédiaires complexes.

\subsubsection{Objectif scientifique et Pipeline d'acquisition des
données}\label{objectif-scientifique-et-pipeline-dacquisition-des-donnuxe9es}

Le métamodèle d'IA vise à apprendre la relation de transfert complexe
entre la morphologie d'une ville, la charge dynamique imposée et ses
conséquences environnementales : \[f : \mathcal{X} \to \mathcal{Y}\] où
: * \(\mathcal{X}\) est l'espace des caractéristiques structurelles et
spectrales d'une ville (43 descripteurs). * \(\mathcal{Y}\) représente
les émissions de \(CO_2\) ou la vitesse moyenne globale générées par la
simulation SUMO de référence (Ground Truth).

L'acquisition et la consolidation du jeu de données d'apprentissage ont
suivi un protocole rigoureux d'automatisation : 1. \textbf{Sélection et
filtrage automatique des simulations} : Un script parcourt de manière
récursive les répertoires d'output. Il trie et ne conserve que les
\textbf{19 simulations complètes et uniques} (les simulations partielles
ou interrompues affichant des émissions nulles ou des fichiers de
métadonnées absents étant écartées). 2. \textbf{Extraction granulaire
des données d'agents} : Les fichiers XML \texttt{tripinfo.xml} issus de
chaque simulation SUMO de référence sont lus à l'aide d'un parseur
événementiel à mémoire constante
(\texttt{xml.etree.ElementTree.iterparse}). Pour chaque agent,
l'algorithme extrait son type de motorisation et sa durée de transit
exacte. 3. \textbf{Découplage fin des taux d'électrification (EV)} :
Pour refléter la réalité des politiques de décarbonation, le modèle
n'utilise pas de taux d'EV global uniforme, mais segmente
l'électrification selon quatre taux catégoriels indépendants : *
\texttt{pct\_car\_ev} (Taux d'électrification des voitures
particulières) * \texttt{pct\_bus\_ev} (Taux d'électrification de la
flotte de bus - ex. VinBus à Hanoï) * \texttt{pct\_truck\_ev} (Taux
d'électrification des camions et livraisons) * \texttt{pct\_moto\_ev}
(Taux d'électrification des deux-roues motorisés)

\subsubsection{Répartition Automatique de la Flotte par Région (Presets
Régionaux)}\label{ruxe9partition-automatique-de-la-flotte-par-ruxe9gion-presets-ruxe9gionaux}

Afin de simplifier l'utilisation du métamodèle et de fiabiliser les
prédictions, le système intègre des distributions de flotte et des taux
d'électrification par défaut calibrés par région géographique d'origine
(Europe, Asie, Afrique, Amérique du Nord, Amérique du Sud, Océanie). Ce
paramétrage automatique sous le capot évite à l'utilisateur d'avoir à
spécifier manuellement des données de répartition complexes, tout en
reflétant fidèlement les disparités géographiques réelles : *
\textbf{Europe \& Amérique du Nord} : Domination des voitures
particulières (\(80\ %
\)) et des véhicules utilitaires, faible proportion de deux-roues
(\(5\ %
\)). Les bus représentent environ \(5\ %
\) de la flotte. * \textbf{Asie du Sud-Est} : Prédominance massive des
deux-roues motorisés (\(60\ %
\)) et des transports collectifs (\(10\ %
\)), avec une dynamique d'électrification rapide des bus urbains portée
par l'écosystème local. * \textbf{Afrique} : Part importante de
véhicules de livraison/camions et de minibus de transport collectif,
avec un taux de pénétration des EV naissant (\(1-2\ %
\)).

\subsubsection{L'architecture prédictive directe à 43
features}\label{larchitecture-pruxe9dictive-directe-uxe0-43-features}

Les émissions de \(CO_2\), qui constituent le but principal de notre
métamodèle d'intelligence artificielle, sont estimées de manière directe
à partir de notre espace de caractéristiques de dimension 43.

Plutôt que d'utiliser un modèle en cascade complexe (qui estimerait
d'abord la vitesse moyenne sur le réseau pour ensuite en déduire les
émissions), notre framework réalise un entraînement direct et optimal :
* \textbf{Modèle d'Émissions de CO2}
(\texttt{xgb\_co2\_predictor.joblib}) : Estime directement la masse
totale d'émissions de \(CO_2\) (en kg) rejetée sur la voirie à partir
des 43 caractéristiques structurelles, géométriques, spectrales et de
charge de trafic.

Cette architecture directe supprime les dépendances intermédiaires,
évitant ainsi la propagation d'erreurs d'estimation d'une étape à
l'autre et garantissant une robustesse accrue.

\begin{quote}
{[}!NOTE{]} \textbf{Généralisation de l'approche à d'autres variables
(ex: la vitesse) :} Bien que la prédiction directe du \(CO_2\) soit le
cœur fonctionnel et opérationnel de notre outil, cette méthodologie de
transfert spectrale est parfaitement transposable à d'autres variables
d'écoulement du trafic. Il est par exemple tout à fait possible de
calquer exactement ce protocole pour entraîner un modèle jumeau
(\texttt{xgb\_speed\_predictor.joblib}) visant à prédire la vitesse
moyenne de transit (\texttt{avg\_speed\_mps}) en exploitant le même
espace d'entrée de 43 caractéristiques. Cela montre la flexibilité et la
portée de la ``signature spectrale'' pour estimer n'importe quelle
métrique macroscopique du réseau.
\end{quote}

Voici la description exhaustive, l'utilité opérationnelle et la
justification scientifique de ces 43 features, classées en 7 catégories
:

\paragraph{Paramètres de Trafic \& Charge (6
features)}\label{paramuxe8tres-de-trafic-charge-6-features}

Ces features mesurent la demande cinématique brute imposée au réseau
urbain lors du scénario d'étude. Elles définissent le terme source de la
congestion. * \textbf{\texttt{nb\_total\_veh} (Nombre total de
véhicules)} : Nombre cumulé de véhicules injectés dans le réseau sur la
durée de la simulation. \emph{Utilité} : C'est le descripteur d'échelle
principal, directement corrélé au volume de émissions de base. *
\textbf{\texttt{duree\_sim\_s} (Durée de simulation, en secondes)} :
Temps d'exposition du réseau à la charge de trafic (typiquement 3600
secondes). \emph{Utilité} : Permet de normaliser les flux temporels et
de distinguer les simulations à court terme des charges continues. *
\textbf{\texttt{nb\_voitures} (Nombre de voitures particulières)} :
Nombre absolu de voitures de tourisme (thermiques et électriques).
\emph{Utilité} : Les voitures particulières forment le socle de la
flotte de base dans la plupart des réseaux. * \textbf{\texttt{nb\_motos}
(Nombre de deux-roues motorisés)} : Nombre absolu de motos et scooters.
\emph{Utilité} : Indispensable pour calibrer le trafic dans les villes
d'Asie du Sud-Est (comme Hanoï) où les deux-roues dominent les flux et
causent du tissage latéral. * \textbf{\texttt{nb\_camions} (Nombre de
camions et véhicules lourds)} : Nombre absolu de poids lourds et
véhicules de livraison. \emph{Utilité} : Les camions ont un facteur
d'émission de CO₂ unitaire très élevé. Même en faible proportion, leur
nombre influence fortement les pics d'émissions. *
\textbf{\texttt{nb\_bus} (Nombre d'autobus)} : Nombre absolu de bus de
transport en commun. \emph{Utilité} : Les bus ont des profils de
conduite spécifiques (arrêts réguliers, accélération lente) qui
influencent la dynamique locale des voies réservées ou partagées.

\paragraph{Taux d'Électrification Découplés (4
features)}\label{taux-duxe9lectrification-duxe9coupluxe9s-4-features}

Ces descripteurs découplent la transition de flotte par catégorie de
véhicules, permettant d'évaluer des politiques environnementales
ciblées. Les véhicules électriques ayant des émissions de CO₂ directes
nulles dans SUMO, ces features agissent comme des modérateurs
d'émissions. * \textbf{\texttt{pct\_car\_ev} (Taux d'électrification des
voitures)} : Pourcentage de voitures électriques (\(0\) à \(100\ %
\)). \emph{Utilité} : Mesure l'impact de l'adoption individuelle de
l'électromobilité. * \textbf{\texttt{pct\_bus\_ev} (Taux
d'électrification de la flotte de bus)} : Pourcentage de bus électriques
(\(0\) à \(100\ %
\)). \emph{Utilité} : Évalue l'impact de l'électrification des flottes
publiques (ex. réseau VinBus). * \textbf{\texttt{pct\_truck\_ev} (Taux
d'électrification des camions)} : Pourcentage de camions électriques
(\(0\) à \(100\ %
\)). \emph{Utilité} : Simule la décarbonation de la logistique urbaine
du dernier kilomètre. * \textbf{\texttt{pct\_moto\_ev} (Taux
d'électrification des deux-roues)} : Pourcentage de deux-roues
électriques (\(0\) à \(100\ %
\)). \emph{Utilité} : Crucial pour mesurer la réduction de pollution
sonore et d'émissions locales dans les métropoles asiatiques saturées de
deux-roues.

\paragraph{Caractéristiques Topologiques Générales (9
features)}\label{caractuxe9ristiques-topologiques-guxe9nuxe9rales-9-features}

Ces features décrivent la structure géométrique et la squelettisation
brute du réseau routier à partir des fichiers XML compilés par
netconvert. * \textbf{\texttt{nodes} (Nœuds)} : Nombre total
d'intersections du réseau (\(|V|\)). \emph{Utilité} : Indique la taille
et la complexité brute de la carte urbaine. * \textbf{\texttt{edges}
(Arêtes)} : Nombre total de segments de voirie orientés (\(|E|\)).
\emph{Utilité} : Mesure la longueur totale et la capacité de stockage
géométrique des voies de la ville. * \textbf{\texttt{density} (Densité
du réseau)} : Ratio du nombre d'arêtes par unité de surface
(\(|E|/\text{Area}\)). \emph{Utilité} : Distingue les réseaux aérés et
suburbains (faible densité) des hyper-centres urbains denses (forte
densité). * \textbf{\texttt{avg\_degree} (Degré moyen)} : Degré de
connectivité moyen des intersections (\(\frac{2|E|}{|V|}\)).
\emph{Utilité} : Un degré moyen élevé (proche de 4) indique des
carrefours à quatre directions, offrant une plus grande flexibilité de
routage. * \textbf{\texttt{asymmetry\_index} (Indice d'asymétrie)} :
Proportion d'arêtes à sens unique dans le réseau
(\(1.0 - |E_{double\_sens}|/|E|\)). \emph{Utilité} : Traduit la
contrainte de circulation directionnelle. Un indice proche de 1 signale
une ville à sens uniques dominants, augmentant la distance de parcours
et le risque de saturation locale. * \textbf{\texttt{sources\_count}
(Nombre de sources)} : Nombre de nœuds n'ayant que des voies sortantes
(degré entrant nul). \emph{Utilité} : Identifie les points d'injection
de trafic périphérique (autoroutes d'accès). *
\textbf{\texttt{sinks\_count} (Nombre de puits)} : Nombre de nœuds
n'ayant que des voies entrantes (degré sortant nul). \emph{Utilité} :
Identifie les zones de sortie ou de stationnement terminal du trafic. *
\textbf{\texttt{sources\_ratio} (Ratio de sources)} : Proportion de
nœuds sources par rapport à la taille du graphe
(\(sources\_count/|V|\)). \emph{Utilité} : Caractérise le profil
d'alimentation du réseau. * \textbf{\texttt{sinks\_ratio} (Ratio de
puits)} : Proportion de nœuds puits par rapport à la taille du graphe
(\(sinks\_count/|V|\)). \emph{Utilité} : Caractérise le profil de
décharge du trafic.

\paragraph{Propriétés Spectrales Non-Pondérées (5
features)}\label{propriuxe9tuxe9s-spectrales-non-ponduxe9ruxe9es-5-features}

Ces descripteurs mathématiques décrivent la stabilité dynamique
intrinsèque du réseau routier modélisé sous forme de graphe orienté
asymétrique, sans tenir compte des longueurs physiques des rues. *
\textbf{\texttt{non\_normalness} (Indice de non-normalité)} : Norme de
Frobenius du commutateur de la matrice d'adjacence non pondérée,
\(\|AA^T - A^TA\|_F\). \emph{Utilité} : C'est le prédicteur des
phénomènes d'amplification transitoire (les ondes de choc et congestions
diffuses inattendues provoquées par une perturbation locale). *
\textbf{\texttt{spectral\_radius} (Rayon spectral, \(\rho\))} : Module
de la valeur propre dominante de la matrice d'adjacence non pondérée.
\emph{Utilité} : Caractérise la capacité globale de transit et la force
des axes structurants (corridors). * \textbf{\texttt{sigma\_max} (Valeur
singulière maximale, \(\sigma_{max}\))} : Norme spectrale induite de la
matrice d'adjacence non pondérée. \emph{Utilité} : Représente le gain
d'amplification maximal théorique du réseau dans le pire scénario de
charge de trafic. * \textbf{\texttt{h2\_norm} (Norme de Hardy \(H_2\))}
: Mesure de l'énergie stockée par la fonction de transfert du réseau
routier. \emph{Utilité} : Quantifie la mémoire de la congestion,
c'est-à-dire le temps nécessaire au réseau pour purger l'excédent de
trafic après un pic d'affluence. * \textbf{\texttt{kreiss\_constant}
(Constante de Kreiss, \(K\))} : Borne supérieure de la résolvante de la
matrice d'adjacence. \emph{Utilité} : Agit comme le « détecteur de
nervosité » du graphe. Une constante de Kreiss élevée préjuge d'un
réseau très sensible aux embouteillages en cascade (\emph{gridlocks}).

\paragraph{Espace Multidimensionnel Top 5 (10
features)}\label{espace-multidimensionnel-top-5-10-features}

Pour capturer la structure hiérarchique complexe et la modularité des
réseaux (les quartiers, les barrières géographiques comme les fleuves,
et la redondance locale), nous avons élargi l'espace spectral en
intégrant les 5 premières valeurs propres et valeurs singulières. *
\textbf{\texttt{lambda\_1}, \texttt{lambda\_2}, \texttt{lambda\_3},
\texttt{lambda\_4}, \texttt{lambda\_5}} : Magnitudes des 5 premières
valeurs propres dominantes ordonnées de la matrice d'adjacence.
\emph{Utilité} : La répartition de ces valeurs propres contient des
informations géométriques sur la structure communautaire du réseau
(spectral gap). Par exemple, un grand écart \(\lambda_1 - \lambda_2\)
indique un réseau homogène, tandis qu'un faible écart révèle des goulots
d'étranglement majeurs (ponts séparant deux zones de la ville). *
\textbf{\texttt{sigma\_1}, \texttt{sigma\_2}, \texttt{sigma\_3},
\texttt{sigma\_4}, \texttt{sigma\_5}} : Les 5 plus grandes valeurs
singulières ordonnées de la matrice d'adjacence. \emph{Utilité} : Si les
premières valeurs singulières sont très proches
(\(\sigma_1 \approx \sigma_2 \approx \sigma_3\)), le réseau dispose de
plusieurs itinéraires alternatifs orthogonaux de capacités similaires
(haute redondance). Si \(\sigma_1 \gg \sigma_2\), le réseau est
monocentrique et vulnérable au blocage complet de son axe unique.

\paragraph{Propriétés Spectrales Pondérées (Physiques \& Capacitives) (3
features)}\label{propriuxe9tuxe9s-spectrales-ponduxe9ruxe9es-physiques-capacitives-3-features}

Ces features intègrent l'impédance géométrique réelle de la voirie (la
longueur de la rue \(L\), sa vitesse autorisée \(W\), et son nombre de
voies \(C\)) au sein de la matrice d'adjacence pondérée
\(A_{ij} = \frac{L_{ij} \cdot C_{ij}}{W_{ij}}\). *
\textbf{\texttt{spectral\_radius\_weighted} (Rayon spectral pondéré)} :
Rayon spectral de la matrice pondérée. \emph{Utilité} : Mesure la
capacité de transport physique réelle du réseau, en tenant compte de la
géométrie capacitive. * \textbf{\texttt{sigma\_max\_weighted} (Valeur
singulière maximale pondérée)} : Plus grande valeur singulière de la
matrice pondérée. \emph{Utilité} : C'est le descripteur spectral le plus
important pour la prédiction de CO₂ (voir Section 6.6), car il lie
l'amplification dynamique aux capacités de stockage et de vitesse
réelles de la voirie. * \textbf{\texttt{h2\_norm\_weighted} (Norme
\(H_2\) pondérée)} : Norme de Hardy \(H_2\) de la matrice pondérée.
\emph{Utilité} : Quantifie la rétention de congestion en secondes
réelles de temps de parcours perdu pour l'ensemble des conducteurs.

\paragraph{Origine Géographique (Encodage One-Hot) (6
features)}\label{origine-guxe9ographique-encodage-one-hot-6-features}

Ces variables catégorielles encodent l'origine continentale du réseau
routier pour servir de proxy aux comportements de conduite locaux et aux
distributions typiques de flotte par défaut : *
\textbf{\texttt{origin\_Europe}} : Caractérise les réseaux radiaux
organiques (ex. Paris, Versailles). *
\textbf{\texttt{origin\_Asia\_Middle\_East}} : Caractérise les mégapoles
denses à trafic mixte (ex. Hanoï, Dubaï). *
\textbf{\texttt{origin\_Africa}} : Caractérise les réseaux en
développement rapide (ex. Nairobi, Casablanca). *
\textbf{\texttt{origin\_Oceania}} : Caractérise les villes côtières
planifiées (ex. Sydney). * \textbf{\texttt{origin\_North\_America}} :
Caractérise les structures en grilles orthogonales régulières (ex. Los
Angeles). * \textbf{\texttt{origin\_South\_America}} : Caractérise les
géométries mixtes côtières (ex. Rio de Janeiro).

\subsubsection{Modèle Mathématique de
XGBoost}\label{moduxe8le-mathuxe9matique-de-xgboost}

L'algorithme XGBoost (eXtreme Gradient Boosting) minimise une fonction
d'apprentissage objective régularisée à l'étape \(t\) pour l'arbre
\(f_t\) :
\[\mathcal{L}^{(t)} = \sum_{i=1}^N l\left(y_i, \hat{y}_i^{(t-1)} + f_t(x_i)\right) + \Omega(f_t)\]
Où le terme de régularisation hybride L1/L2 est défini par :
\[\Omega(f_t) = \gamma T + \frac{1}{2} \lambda \sum_{j=1}^T w_j^2 + \alpha \sum_{j=1}^T |w_j|\]
XGBoost résout ce problème en appliquant un développement de Taylor au
second ordre de la perte :
\[\mathcal{L}^{(t)} \approx \sum_{i=1}^N \left[ g_i f_t(x_i) + \frac{1}{2} h_i f_t^2(x_i) \right] + \gamma T + \frac{1}{2} \lambda \sum_{j=1}^T w_j^2\]
où \(g_i\) (gradient) et \(h_i\) (hessienne) désignent les dérivées
première et seconde de la perte par rapport à la prédiction précédente :
\[g_i = \frac{\partial l\left(y_i, \hat{y}_i^{(t-1)}\right)}{\partial \hat{y}_i^{(t-1)}} \quad \text{et} \quad h_i = \frac{\partial^2 l\left(y_i, \hat{y}_i^{(t-1)}\right)}{\partial \left(\hat{y}_i^{(t-1)}\right)^2}\]
Cette formulation permet d'intégrer de manière optimale les gradients de
second ordre, ce qui confère au modèle une vitesse de convergence
exceptionnelle et une grande stabilité face aux interactions non
linéaires.

\subsubsection{Analyse comparative des performances réelles
(CO2)}\label{analyse-comparative-des-performances-ruxe9elles-co2}

Le métamodèle direct a été entraîné sur la base consolidée des 19
simulations d'échelle réaliste. Les performances obtenues en validation
croisée pour la prédiction directe du \(CO_2\) à partir des 43 features
sont résumées ci-dessous :

{\def\LTcaptype{none} % do not increment counter
\begin{longtable}[]{@{}
  >{\raggedright\arraybackslash}p{(\linewidth - 6\tabcolsep) * \real{0.2105}}
  >{\centering\arraybackslash}p{(\linewidth - 6\tabcolsep) * \real{0.2632}}
  >{\centering\arraybackslash}p{(\linewidth - 6\tabcolsep) * \real{0.2632}}
  >{\centering\arraybackslash}p{(\linewidth - 6\tabcolsep) * \real{0.2632}}@{}}
\toprule\noalign{}
\begin{minipage}[b]{\linewidth}\raggedright
Modèle
\end{minipage} & \begin{minipage}[b]{\linewidth}\centering
RMSE (\(CO_2\) en kg)
\end{minipage} & \begin{minipage}[b]{\linewidth}\centering
MAE (\(CO_2\) en kg)
\end{minipage} & \begin{minipage}[b]{\linewidth}\centering
Score \(R^2\) (\%)
\end{minipage} \\
\midrule\noalign{}
\endhead
\bottomrule\noalign{}
\endlastfoot
\textbf{XGBoost (Direct)} & \textbf{15 888,62} & \textbf{13 343,48} &
\textbf{95,03 \%} \\
\textbf{Régression Ridge (L2)} & 30 255,52 & 28 579,07 & 81,99 \% \\
\textbf{Random Forest} & 32 239,28 & 30 526,94 & 79,56 \% \\
\textbf{MLP Regressor (Réseau de Neurones)} & 194 716,32 & 141 992,69 &
-645,78 \% \\
\end{longtable}
}

\paragraph{Analyse critique des résultats
:}\label{analyse-critique-des-ruxe9sultats}

\begin{itemize}
\tightlist
\item
  \textbf{XGBoost (Direct)} : Il surclasse largement les autres
  approches avec un \(R^2\) de \(95,03\ \%\). La prédiction directe à
  partir des descripteurs spectraux et physiques permet de capter la
  non-linéarité des émissions carbone sans avoir besoin d'ajustements
  intermédiaires.
\item
  \textbf{Random Forest} : Bien que robuste, sa nature de moyennage
  (bagging) atténue les pics non linéaires de pollution survenant lors
  de congestions géométriques brutales, limitant son score à
  \(79,56\ \%\).
\item
  \textbf{MLP Regressor} : Subit un échec complet en validation croisée
  (\(R^2 = -645,78\ %
  \)). Ce comportement s'explique par la taille limitée de notre jeu de
  données (19 points de simulations complètes à grande échelle). Les
  réseaux de neurones profonds exigent des milliers d'échantillons pour
  converger sans sur-apprentissage massif.
\end{itemize}

\subsubsection{Importance des variables (Feature Importance
CO2)}\label{importance-des-variables-feature-importance-co2}

L'extraction de l'importance relative des descripteurs dans le modèle
XGBoost de prédiction du \(CO_2\) met en lumière la hiérarchie
d'influence physique et spectrale suivante :

\begin{enumerate}
\def\labelenumi{\arabic{enumi}.}
\tightlist
\item
  \textbf{\texttt{nb\_total\_veh} (31,64 \%)} : Le volume global de
  circulation reste le premier facteur d'effort du réseau.
\item
  \textbf{\texttt{nb\_voitures} (17,24 \%)} : La flotte automobile de
  tourisme, principal contributeur de base.
\item
  \textbf{\texttt{nb\_motos} (16,60 \%)} : Les deux-roues motorisés,
  caractéristique des morphologies d'Asie du Sud-Est.
\item
  \textbf{\texttt{edges} (14,51 \%)} : Nombre de rues de la ville
  (facteur dimensionnel topologique).
\item
  \textbf{\texttt{nb\_camions} (12,22 \%)} : Volume de camions lourds à
  fort taux unitaire de rejet.
\item
  \textbf{\texttt{duree\_sim\_s} (4,34 \%)} : Temps total d'exposition
  et d'écoulement du trafic.
\item
  \textbf{\texttt{nb\_bus} (1,70 \%)} : Part des bus de transport
  public.
\item
  \textbf{\texttt{pct\_moto\_ev} (0,87 \%)} : Pénétration électrique des
  deux-roues.
\item
  \textbf{\texttt{sigma\_max\_weighted} (0,37 \%)} : Plus grande valeur
  singulière de la matrice d'adjacence pondérée (capacité physique).
\item
  \textbf{\texttt{origin\_Europe} (0,16 \%)} : Correction géographique
  des flux.
\end{enumerate}

Ces résultats valident scientifiquement l'apport des métriques
spectrales pondérées comme ajusteurs de congestion dans les modèles de
machine learning.

\subsubsection{Espace de transfert topologique : Généralisation
Cross-City}\label{espace-de-transfert-topologique-guxe9nuxe9ralisation-cross-city}

Le défi ultime de ce moteur prédictif réside dans sa capacité de
généralisation transversale (\emph{cross-city generalization}) : estimer
instantanément les performances d'une ville sans aucune simulation SUMO
préalable.

Pour ce faire, nous formalisons l'apprentissage par transfert sous la
forme d'une interpolation barycentrique au sein de l'espace des
signatures spectrales. Soit \(\vec{x}_{target}\) le vecteur des
descripteurs spectraux et topologiques normalisés d'une ville cible (ex.
Nairobi), extrait instantanément à partir de ses données OpenStreetMap.

Nous exprimons ce vecteur comme une combinaison linéaire convexe des
vecteurs de descripteurs de \(k\) villes déjà simulées et présentes dans
la base d'entraînement :
\[\vec{x}_{target} = \sum_{j=1}^k \alpha_j \vec{x}_j\] sous les
contraintes de fermeture convexe :
\[\sum_{j=1}^k \alpha_j = 1 \quad \text{et} \quad \alpha_j \ge 0 \quad \forall j \in \{1, \dots, k\}\]

Les coefficients de pondération \(\alpha_j\) (les coordonnées
barycentriques spectrales) sont calculés en résolvant un problème
d'optimisation quadratique de distance minimale :
\[\vec{\alpha}^* = \arg\min_{\vec{\alpha}} \left\| \vec{x}_{target} - \sum_{j=1}^k \alpha_j \vec{x}_j \right\|_2^2\]

Si le profil spectral de Nairobi se révèle géométriquement proche de
Marseille (poids \(\alpha_{Marseille} = 0.70\)) et de Berlin (poids
\(\alpha_{Berlin} = 0.30\)), l'intelligence artificielle estime la
courbe d'émissions de la cible en exploitant les surfaces de décision
apprises sur ces analogues morphologiques. Cette généralisation
spectrale (illustrée par les diagrammes de dispersion et courbes de
convergence en Annexe D) ramène le temps d'évaluation d'un nouveau plan
urbain de plusieurs heures de micro-simulation à quelques millisecondes
d'inférence.

\newpage

\section{TROISIÈME PARTIE : APPLICATIONS EXPÉRIMENTALES ET VALIDATION
COMPARATIVE}\label{troisiuxe8me-partie-applications-expuxe9rimentales-et-validation-comparative}

\subsection{Chapitre 8 : Étude de cas localisée -- Le hub de recharge
rapide de Sao Bien
(Hanoï)}\label{chapitre-8-uxe9tude-de-cas-localisuxe9e-le-hub-de-recharge-rapide-de-sao-bien-hanouxef}

\subsubsection{Morphologie géométrique et contraintes d'infrastructure
du site
d'étude}\label{morphologie-guxe9omuxe9trique-et-contraintes-dinfrastructure-du-site-duxe9tude}

Le complexe résidentiel et commercial de Vinhomes Ocean Park (VHOP),
situé à la périphérie Est de la métropole de Hanoï, constitue une
opportunité d'étude unique pour analyser la pénétration à grande échelle
de l'électromobilité en milieu urbain hyper-dense. Notre recherche se
focalise spécifiquement sur le corridor de l'avenue Sao Bien, un axe
routier bidirectionnel structurant qui dessert l'entrée principale du
centre commercial \emph{Vincom Mega Mall}.

Le site d'étude concentre deux infrastructures de service contiguës
situées sur une voie latérale de desserte commerciale à faible capacité
(largeur de voie de 3,5 mètres, limitée à deux voies de circulation
unidirectionnelle après insertion depuis l'avenue principale) : 1.
\textbf{Une zone d'attente de taxis (GSM Taxi Waiting Zone) :}
Positionnée immédiatement en amont du centre commercial, disposant de 16
places de stockage organisées en stationnement parallèle le long de la
chaussée. 2. \textbf{Le Hub de recharge ultra-rapide VinFast (V-Green
Super-fast Charging Hub) :} Situé à environ 75 mètres après
l'intersection d'entrée, configuré avec 12 emplacements de recharge en
épi à 135 degrés par rapport au sens de circulation de la voie de
service. Cette géométrie spécifique impose aux véhicules un recul à 90
degrés pour s'extraire de la borne de recharge, créant un blocage
temporaire complet de la voie de circulation adjacente lors de chaque
manœuvre de départ.

\begin{verbatim}
                                  Avenue Sao Bien (Axe Principal)
   ========================================================================================
             ||  Insertion
             \/
   ----------------------------------------------------------------------------------------
   [Voie de Dessers]  ==>  [Taxi Waiting Zone (16 slots)]  ==>  [EV Charging Hub (12 slots)]
   ----------------------------------------------------------------------------------------
\end{verbatim}

La problématique physique du site réside dans la superposition de ces
usages au sein d'une géométrie contrainte. Les taxis électriques GSM en
attente de prise en charge et les véhicules électriques (particuliers et
taxis) convergeant vers les bornes de recharge partagent la même voie
d'accès. De plus, la forte prédominance des motocyclistes dans le flux
de trafic de Hanoï (qui représente entre 50 \% et 64 \% de la
composition globale) engendre un phénomène de ``tissage'' permanent. Les
deux-roues s'insèrent dans les micro-intervalles géométriques laissés
par les voitures en manœuvre, augmentant la friction cinématique et
ralentissant les manœuvres d'insertion et d'extraction du hub de
recharge.

\subsubsection{Modélisation cinématique et calibrage stochastique de la
recharge}\label{moduxe9lisation-cinuxe9matique-et-calibrage-stochastique-de-la-recharge}

\paragraph{Le modèle de temps d'arrêt stochastique (Stochastic Dwell
Time)}\label{le-moduxe8le-de-temps-darruxeat-stochastique-stochastic-dwell-time}

Pour représenter fidèlement l'impact de la station de recharge sur le
trafic, il est scientifiquement inexact de modéliser le temps de
raccordement des véhicules électriques (\emph{dwell time}) par une
valeur fixe ou une moyenne simpliste. Dans la réalité physique, la durée
d'une session de recharge est une variable stochastique complexe. Elle
dépend de la puissance nominale de la borne (150 kW DC), mais également
de la courbe de charge de la batterie régie par le système de gestion
thermique (Battery Management System - BMS), du niveau de charge initial
du véhicule à son arrivée (\(SoC_{in}\)) et du niveau de charge souhaité
par l'usager à son départ (\(SoC_{out}\)).

Selon les spécifications techniques de VinFast et de l'opérateur de
recharge V-Green, le cycle de charge standard pour une batterie de 42 à
80 kWh (modèles VF e34, VF8, VF9) de 10 \% à 70 \% de \(SoC\) s'établit
à \textbf{31 minutes} sous une puissance stabilisée. Au-delà de 70 \%,
la puissance de charge décroît de manière exponentielle pour préserver
l'intégrité électrochimique des cellules.

Pour intégrer cette variabilité au sein du simulateur SUMO, nous avons
développé un \textbf{Modèle de Temps d'Arrêt Stochastique} basé sur une
loi normale tronquée. La durée d'arrêt \(d\) d'un véhicule électrique
ciblé pour une session de recharge suit la loi :
\[d \sim \mathcal{N}(\mu, \sigma^2)\] calibrée selon quatre profils
opérationnels d'encombrement du hub issus de nos observations : *
\textbf{Profil Léger (Light) :} Modélise des recharges d'appoint
rapides.
\[\mu = 1050 \text{ s } (17.5 \text{ min}), \quad \sigma = 225 \text{ s } (3.75 \text{ min})\]
* \textbf{Profil Modéré (Moderate) :} Calibré sur le benchmark standard
VinFast.
\[\mu = 2100 \text{ s } (35 \text{ min}), \quad \sigma = 300 \text{ s } (5 \text{ min})\]
* \textbf{Profil Lourd (Heavy) :} Modélise les sessions de recharge
complètes ou les ralentissements dus au partage de puissance dynamique
(\emph{power-sharing}) entre bornes adjacentes.
\[\mu = 3600 \text{ s } (60 \text{ min}), \quad \sigma = 450 \text{ s } (7.5 \text{ min})\]
* \textbf{Profil Critique (Critical) :} Représente les situations de
saturation sévère et le comportement d'indiscipline des usagers
(``overstaying''), où les véhicules restent connectés ou stationnés
après la fin effective de leur cycle de charge.
\[\mu = 5400 \text{ s } (90 \text{ min}), \quad \sigma = 900 \text{ s } (15 \text{ min})\]

\paragraph{Le mécanisme d'initialisation dynamique par
``Warm-Start''}\label{le-muxe9canisme-dinitialisation-dynamique-par-warm-start}

Les simulateurs de trafic microscopiques démarrent classiquement dans un
état ``vide et froid'' (\emph{cold-start}), où aucun véhicule n'est
présent sur le réseau à l'instant \(t = 0\). Pour une modélisation de
flux continu, cette approche introduit un biais transitoire important :
durant les 15 à 20 premières minutes de la simulation, le hub de
recharge est vide, ce qui fausse les temps d'attente et sous-estime la
congestion d'accès.

Pour résoudre ce verrou de fidélité, nous avons conçu un algorithme
d'initialisation dynamique par \textbf{``Warm-Start'' (démarrage à
chaud)}. À l'instant \(t=0\), avant l'injection du premier véhicule
actif de la simulation, le hub est pré-peuplé stochastiquement avec des
\textbf{véhicules fantômes} (\emph{ghost vehicles}) occupant un nombre
prédéterminé de bornes de recharge.

Le nombre de bornes occupées à l'initialisation est calibré selon le
niveau de charge du scénario : * \emph{Scénarios fluides / creux
(Midday) :} 7 à 9 slots occupés stochastiquement (taux d'occupation
initial de 58 \% à 75 \%). * \emph{Scénarios de pointe (Rush Hour /
Holiday) :} 11 à 12 slots occupés (taux d'occupation initial de 91 \% à
100 \%).

Pour éviter un départ simultané de ces véhicules qui créerait une onde
de choc artificielle, l'algorithme attribue à chaque véhicule fantôme
\(j\) un \textbf{temps d'occupation résiduel \(t_{res, j}\)} tiré
aléatoirement selon une distribution uniforme :
\[t_{res, j} \sim \mathcal{U}(300, 900) \text{ secondes}\]

Ainsi, dès la première seconde de la simulation physique, les bornes de
recharge présentent un niveau d'encombrement réaliste et libèrent leurs
places de manière échelonnée dans le temps, offrant une reproduction
fidèle des dynamiques d'attente et d'insertion vécues par les
conducteurs arrivant sur site.

\subsubsection{Analyse des données réelles et identification des profils
empiriques}\label{analyse-des-donnuxe9es-ruxe9elles-et-identification-des-profils-empiriques}

\paragraph{Constitution du jeu de données et exclusion systématique des
bus}\label{constitution-du-jeu-de-donnuxe9es-et-exclusion-systuxe9matique-des-bus}

Pour calibrer le jumeau numérique, nous exploitons les comptages de
trafic issus de notre campagne de mesures vidéo par vision artificielle
(YOLOv8) sur l'avenue Sao Bien. Le jeu de données consolidé comporte 72
enregistrements unitaires de 10 minutes, couvrant différentes journées
et tranches horaires entre le 14 avril et le 7 mai 2026.

Dans notre formalisation des flux cinématiques et environnementaux,
\textbf{les bus de transport en commun ont été systématiquement exclus
des statistiques de composition et des temps d'attente du hub}. Cette
décision méthodologique s'appuie sur deux fondements scientifiques : 1.
\textbf{Indépendance fonctionnelle :} Les bus électriques circulants
(flotte \emph{VinBus}) opèrent sur des itinéraires fixes et rigides avec
des fréquences déterminées (2.0 à 3.0 passages par intervalle de 10
minutes en moyenne). Ils ne s'insèrent jamais dans le hub de recharge
des véhicules légers, disposant de leurs propres stations de charge
dédiées au dépôt principal. 2. \textbf{Distorsion statistique :}
Intégrer les bus dans les calculs de répartition globale sous-estimerait
l'importance relative de la flotte de voitures particulières et de
taxis, qui constitue 100 \% de la demande physique pesant sur le hub de
recharge de Sao Bien.

\paragraph{Caractérisation des trois profils
empiriques}\label{caractuxe9risation-des-trois-profils-empiriques}

Le traitement statistique des données réelles a permis d'isoler trois
profils de trafic distincts et hautement contrastés.

\subparagraph{Le Profil de Ligne de Base (Regular Midday
Baseline)}\label{le-profil-de-ligne-de-base-regular-midday-baseline}

Ce profil, calculé à partir de 58 observations stables lors des heures
creuses de la mi-journée (11h00 - 13h00) hors jours fériés, caractérise
un écoulement fluide à volume modéré : * \textbf{Volume moyen total
(excluant les bus) :} \textbf{134,10 véhicules par tranche de 10
minutes}. * \textbf{Répartition modale :} * \emph{Motorcycles
(deux-roues) :} \textbf{67,97 unités} (\textbf{50,7 \%}). Le mode
deux-roues reste majoritaire, reflétant la structure classique des
déplacements quotidiens à Hanoï. * \emph{Standard Cars (ICE) :}
\textbf{30,38 unités} (\textbf{22,7 \%}). * \emph{Electric Vehicles (EV)
:} \textbf{35,74 unités} (\textbf{26,7 \%}). Ce taux d'électrification
de près de 27 \% de la flotte totale s'explique par la forte
concentration de taxis GSM sur la zone commerciale.

\subparagraph{Le Profil d'Heure de Pointe (Regular Evening
Peak)}\label{le-profil-dheure-de-pointe-regular-evening-peak}

Ce profil modélise la surpression cinématique observée en fin de journée
(17h00 - 18h00), marquée par le retour des travailleurs et la
convergence vers les pôles de restauration du centre commercial : *
\textbf{Volume moyen total (excluant les bus) :} \textbf{227,67
véhicules par tranche de 10 minutes} (soit une augmentation de
\textbf{70 \%} du flux par rapport à la baseline). * \textbf{Répartition
modale :} * \emph{Motorcycles (deux-roues) :} \textbf{143,06 unités}
(\textbf{62,8 \%}). La part des deux-roues s'accroît significativement,
exacerbant le phénomène de friction latérale. * \emph{Standard Cars
(ICE) :} \textbf{34,44 unités} (\textbf{15,1 \%}). * \emph{Electric
Vehicles (EV) :} \textbf{50,17 unités} (\textbf{22,0 \%}).

\subparagraph{Le Profil de Rupture : Le phénomène de ``Holiday
Reversal''}\label{le-profil-de-rupture-le-phuxe9nomuxe8ne-de-holiday-reversal}

Les observations menées lors des journées nationales chômées du 30 avril
et du 1er mai ont révélé une anomalie comportementale majeure.
Contrairement aux modèles classiques qui prévoient une simple homothétie
des volumes, nous avons observé une \textbf{inversion modale
structurelle complète (Holiday Reversal)}.

Le volume de trafic total global diminue légèrement pour atteindre
\textbf{117,17 véhicules par tranche de 10 minutes} en moyenne (en
raison du départ de la population de la ville), mais la composition de
la flotte subit une mutation drastique : * \emph{Motorcycles
(deux-roues) :} \textbf{46,50 unités} (\textbf{39,7 \%}). On observe une
baisse de plus de 10 points de pourcentage de la part des motos par
rapport au midi normal. * \emph{Standard Cars (ICE) :} \textbf{23,83
unités} (\textbf{20,3 \%}). * \emph{Electric Vehicles (EV) :}
\textbf{46,83 unités} (\textbf{40,0 \%}).

Lors de cette période festive, la part des véhicules électriques grimpe
à 40 \% de la flotte totale en circulation. Si l'on restreint l'analyse
aux seuls véhicules à quatre roues (ICE + EV), \textbf{le taux de
pénétration des véhicules électriques atteint le niveau record de 66,3
\%}. Ce phénomène s'explique par l'utilisation massive de la flotte de
taxis GSM par les familles vietnamiennes visitant le centre commercial
et les lagunes artificielles de VHOP, combinée à l'usage des SUV
électriques personnels.

Pour le hub de recharge, ce profil représente un cas d'étude critique de
saturation instantanée, les bornes fonctionnant en continu à 100 \% de
leur capacité nominale avec l'apparition de files d'attente bloquantes
sur la voirie.

\begin{verbatim}
       RÉPARTITION DU FLUIDE MIDDAY (REGULIER VS HOLIDAY REVERSAL)
  ========================================================================
  [Midi Régulier]  |====== Moto (50.7%) =====|=== EV (26.7%) ===|== ICE (22.7%) ==|
  
  [Midi Vacances]  |==== Moto (39.7%) ===|====== EV (40.0%) ======|== ICE (20.3%) ==|
  ========================================================================
\end{verbatim}

\subparagraph{La Période de Transition Pré-Vacances (27 - 29
Avril)}\label{la-puxe9riode-de-transition-pruxe9-vacances-27---29-avril}

Les mesures enregistrées durant les trois jours précédant les fêtes (27,
28 et 29 avril) décrivent un régime transitoire de montée en charge. Le
volume moyen s'établit à \textbf{123,83 véhicules par tranche de 10
minutes}, avec une répartition intermédiaire : * \emph{Motorcycles :}
\textbf{58,00 unités} (\textbf{46,8 \%}). * \emph{Standard Cars (ICE) :}
\textbf{32,33 unités} (\textbf{26,1 \%}). * \emph{Electric Vehicles :}
\textbf{33,50 unités} (\textbf{27,1 \%}).

Cette phase de transition montre que l'inversion modale n'est pas un
événement instantané mais un processus graduel de modification des
usages de mobilité, fournissant des indicateurs précurseurs
indispensables pour anticiper la congestion du réseau d'alimentation
électrique.

\subsubsection{Évaluation des scénarios d'évolution (S0 - S5) et aide à
la
décision}\label{uxe9valuation-des-scuxe9narios-duxe9volution-s0---s5-et-aide-uxe0-la-duxe9cision}

Pour projeter l'impact de la croissance de Vinhomes Ocean Park d'ici
2030, nous avons simulé cinq scénarios d'évolution sur notre jumeau
numérique SUMO :

\paragraph{Scénario S0 \& S1 : Baseline et Heure de pointe historiques
(2023)}\label{scuxe9nario-s0-s1-baseline-et-heure-de-pointe-historiques-2023}

Calibrés directement sur les débits réels observés. Le réseau est stable
dans le scénario S0 (temps d'attente moyen au hub inférieur à 45
secondes pour les EVs). Dans le scénario S1, des congestions localisées
apparaissent sur la voie de service en raison des conflits d'insertion
avec la file d'attente des taxis, mais la fluidité globale sur l'avenue
Sao Bien reste préservée (vitesse moyenne de 32 km/h).

\paragraph{Scénario S2 : Le pic
Holiday}\label{scuxe9nario-s2-le-pic-holiday}

Modélise la configuration du ``Holiday Reversal'' (taux d'adoption EV à
75 \% au sein des véhicules de tourisme). La station de recharge est
saturée dès les premières minutes. Les véhicules électriques en attente
de bornes libres forment une file d'attente de 4 à 6 véhicules sur la
voie de service, provoquant une baisse de la vitesse moyenne sur cet axe
à moins de 12 km/h.

\paragraph{Scénarios S3a \& S3b : Horizon 2028 (60 \% d'occupation
urbaine)}\label{scuxe9narios-s3a-s3b-horizon-2028-60-doccupation-urbaine}

Ces scénarios projettent une augmentation globale des volumes de trafic
(x2.0 pour le midi S3a, x3.2 pour l'heure de pointe S3b) et un taux
d'adoption EV de 80 \% à 85 \%. * \textbf{Résultats :} Dans le scénario
S3a, la file d'attente au hub commence à déborder de la voie de service
pour impacter la voie lente de l'avenue Sao Bien. * Dans le scénario de
pointe S3b, le goulot d'étranglement devient structurel : le temps
d'attente moyen pour un véhicule électrique avant d'accéder à une borne
s'élève à \textbf{18,4 minutes}. La congestion sur l'avenue Sao Bien
remonte sur plus de 150 mètres en amont de l'intersection d'entrée.

\paragraph{Scénarios S4a \& S4b : Horizon 2030 (Saturation à 80 \%
d'occupation)}\label{scuxe9narios-s4a-s4b-horizon-2030-saturation-uxe0-80-doccupation}

Ce scénario modélise le stress-test limite avec une multiplication par
4.0 du volume de trafic global d'heure de pointe et une transition à 100
\% vers l'électromobilité pour les véhicules légers. * \textbf{Résultats
(S4b) :} Le système atteint son point de rupture physique complète. La
voie de service est entièrement paralysée par une file d'attente
ininterrompue de taxis et de véhicules électriques. Les manœuvres
d'extraction des bornes (recul à 135 degrés) bloquent les flux durant
plusieurs minutes. Ce blocage local se propage par effet cascade
(gridlock) sur l'intégralité du carrefour principal de l'avenue Sao
Bien. Le temps d'attente moyen au hub s'envole à \textbf{45,2 minutes},
et la vitesse moyenne sur le réseau s'effondre à \textbf{4,2 km/h}
(vitesse équivalente à de la marche à pied), confirmant la paralysie
fonctionnelle de l'infrastructure.

\paragraph{Scénario S5 : La politique de mitigation par expansion
physique}\label{scuxe9nario-s5-la-politique-de-mitigation-par-expansion-physique}

Face au constat de rupture du scénario S4b, nous avons testé l'impact
d'une décision d'aménagement urbain : l'expansion du hub de recharge à
\textbf{20 slots (+8 bornes de 150 kW DC)} et la création d'une voie
d'insertion dédiée pour la zone d'attente des taxis (voir Figure D.4,
Annexe D). * \textbf{Résultats :} L'évaluation sur le jumeau numérique
SUMO démontre l'efficacité de cette politique de mitigation.
L'augmentation de la capacité de traitement résorbe la file d'attente
sur la voie de service. Le temps d'attente moyen pour accéder à une
borne de recharge chute de 45,2 minutes à \textbf{6,8 minutes} (-85 \%).
Les manœuvres de sortie en épi, bien que toujours gênantes, ne bloquent
plus le flux principal grâce à la voie d'évitement supplémentaire. La
vitesse moyenne globale sur le corridor Sao Bien remonte à \textbf{21,5
km/h}, rétablissant un niveau de service (Level of Service - LOS)
acceptable pour la communauté.

Ce cas d'étude démontre de manière empirique l'utilité du jumeau
numérique microscopique comme outil d'aide à la décision publique pour
planifier la transition énergétique des réseaux de transport urbain.

\newpage

\subsection{Chapitre 9 : Cadre d'expérimentation global et comparative
des réseaux
urbains}\label{chapitre-9-cadre-dexpuxe9rimentation-global-et-comparative-des-ruxe9seaux-urbains}

\subsubsection{Conception d'un framework de simulation multi-ville
agnostique}\label{conception-dun-framework-de-simulation-multi-ville-agnostique}

Pour valider l'indépendance de notre moteur de simulation et sa capacité
à traiter des réseaux de complexités variées, nous avons déployé un
protocole expérimental global impliquant six réseaux urbains aux
caractéristiques contrastées : * \textbf{Versailles :} Réseau aéré de
taille modérée (2 308 nœuds, 9 095 arêtes) caractérisé par une structure
géométrique planifiée. * \textbf{Paris :} Réseau historique de grande
taille (15 735 nœuds, 62 378 arêtes) présentant une structure radiale
dense. * \textbf{Madrid :} Réseau métropolitain étendu (61 258 nœuds,
214 400 arêtes) mêlant un centre dense et des autoroutes urbaines. *
\textbf{Berlin :} Graphe urbain de grande dimension (61 135 nœuds, 213
900 arêtes) caractérisé par une structure mixte. * \textbf{Los Angeles
:} Réseau autoroutier et urbain de très grande taille (176 905 nœuds,
619 100 arêtes) structuré sous forme de grille orthogonale. *
\textbf{Hanoï :} Réseau géant (473 209 nœuds, 1 656 200 arêtes)
modélisant une mégapole asiatique à forte densité et topologie
irrégulière.

Le framework est conçu pour exécuter la même logique de simulation
(chargement du réseau, génération de trajets avec contrainte de distance
minimale, exécution physique de SUMO, et extraction des métriques
écologiques HBEFA3) de manière unifiée, en adaptant uniquement les
chemins d'accès système aux fichiers de topologie respectifs.

\subsubsection{Analyse comparative des résiliences géométriques : Radial
européen vs Grille orthogonale vs Boucles
fermées}\label{analyse-comparative-des-ruxe9siliences-guxe9omuxe9triques-radial-europuxe9en-vs-grille-orthogonale-vs-boucles-fermuxe9es}

L'évaluation des quatre scénarios comportementaux (Constant, Rush Hour,
Max Jam, Bottleneck) sur ces six réseaux révèle l'influence déterminante
de la géométrie de la voirie sur la résilience globale du trafic.

\paragraph{Morphologie Radiale Européenne (Paris,
Madrid)}\label{morphologie-radiale-europuxe9enne-paris-madrid}

Ces réseaux se caractérisent par une forte convergence des axes
principaux vers des points centraux (structures en étoile). Lors du
scénario \emph{Bottleneck}, ces structures se révèlent vulnérables : la
saturation d'un axe central remonte rapidement le long des voies d'accès
radiales, bloquant les carrefours en amont. En raison de la densité des
nœuds historiques et du manque de voies rapides de dérivation, les flux
de trafic ne disposent pas d'alternatives viables, provoquant un
effondrement de la vitesse moyenne et une hausse rapide des émissions de
\(CO_2\).

\paragraph{Grille Orthogonale Nord-Américaine (Los
Angeles)}\label{grille-orthogonale-nord-amuxe9ricaine-los-angeles}

La structure régulière de Los Angeles présente un comportement
différent. Grâce à la régularité des carrefours et à la redondance des
itinéraires parallèles, le réseau absorbe mieux les surcharges locales
du scénario \emph{Bottleneck}. Les véhicules se répartissent d'eux-mêmes
sur les axes adjacents via l'algorithme de routage dynamique. Cependant,
cette structure est sensible au scénario \emph{Max Jam} : le volume
important d'intersections régulées par des feux de signalisation crée
des files d'attente successives qui saturent les intersections si
l'injection de véhicules est trop massive.

\paragraph{Structures Fermées et Voies de Service (Sao Bien,
Hanoï)}\label{structures-fermuxe9es-et-voies-de-service-sao-bien-hanouxef}

La morphologie de Hanoï et des complexes de type Vinhomes se caractérise
par des axes sinueux, des voies de service étroites et des contraintes
d'accès (U-turns imposés par des terre-pleins centraux). Cette géométrie
présente une faible résilience face aux variations de charge. La moindre
obstruction locale (par exemple, un véhicule électrique manœuvrant pour
entrer dans une borne de recharge) bloque l'une des deux voies de
circulation disponibles, forçant les motos à se faufiler et ralentissant
l'ensemble du flux. Ces réseaux présentent un comportement non-linéaire
prononcé, passant sans transition d'un état fluide à une congestion
complète.

\subsubsection{Limites matérielles de la micro-simulation : Surcharge
mémoire (RAM), traitement XML et phénomène de
SWAP}\label{limites-matuxe9rielles-de-la-micro-simulation-surcharge-muxe9moire-ram-traitement-xml-et-phuxe9nomuxe8ne-de-swap}

Les expérimentations menées sur ces six villes ont mis en évidence les
limites physiques des systèmes de calcul pour exécuter des jumeaux
numériques microscopiques sur de grands réseaux
(EXPLICATION\_PERFORMANCE.md).

Le tableau suivant résume la taille des fichiers de topologie sur le
disque dur et l'occupation mémoire vive (RAM) correspondante mesurée en
Python lors de l'exécution de l'étape de routage (chargement du réseau
via \texttt{sumolib}) :

{\def\LTcaptype{none} % do not increment counter
\begin{longtable}[]{@{}
  >{\raggedright\arraybackslash}p{(\linewidth - 6\tabcolsep) * \real{0.2105}}
  >{\centering\arraybackslash}p{(\linewidth - 6\tabcolsep) * \real{0.2632}}
  >{\centering\arraybackslash}p{(\linewidth - 6\tabcolsep) * \real{0.2632}}
  >{\centering\arraybackslash}p{(\linewidth - 6\tabcolsep) * \real{0.2632}}@{}}
\toprule\noalign{}
\begin{minipage}[b]{\linewidth}\raggedright
Ville
\end{minipage} & \begin{minipage}[b]{\linewidth}\centering
Taille du fichier net.xml (Disque)
\end{minipage} & \begin{minipage}[b]{\linewidth}\centering
Occupation mémoire RAM (Python DOM)
\end{minipage} & \begin{minipage}[b]{\linewidth}\centering
Temps de chargement initial
\end{minipage} \\
\midrule\noalign{}
\endhead
\bottomrule\noalign{}
\endlastfoot
\textbf{Versailles} & 6,6 Mo & \textasciitilde110 Mo & \textless{} 1
s \\
\textbf{Paris} & 52 Mo & \textasciitilde880 Mo & 4 s \\
\textbf{Madrid} & 184 Mo & \textasciitilde3,1 Go & 18 s \\
\textbf{Los Angeles} & 533 Mo & \textasciitilde8,9 Go & 55 s \\
\textbf{Hanoï} & 1 239 Mo & \textasciitilde18,5 Go & 142 s \\
\end{longtable}
}

Pour les réseaux de Versailles, Paris et Madrid, l'occupation mémoire
reste compatible avec les architectures de PC de bureau standards
(généralement équipés de 16 Go de RAM).

Cependant, pour le réseau de Hanoï (1,2 Go de données XML sur disque),
l'empreinte mémoire atteint \textbf{18,5 Go de RAM}. Sur une machine
disposant de 16 Go de mémoire physique, ce chargement provoque une
saturation immédiate. Pour maintenir le processus actif, le système
d'exploitation Windows active le mécanisme de \textbf{pagination
virtuelle (SWAP)}, transférant les pages mémoire excédentaires vers le
disque de stockage (SSD ou HDD).

Ce recours au stockage de masse ralentit considérablement la vitesse de
calcul. Le processeur passe une partie de ses cycles d'horloge en
attente d'entrées/sorties disque, allongeant le temps de génération des
itinéraires de quelques secondes à plusieurs dizaines de minutes. Cette
limite matérielle montre la difficulté de déployer des simulations
physiques en temps réel pour l'aide à la décision à l'échelle
métropolitaine (voir Figures D.10 et D.11, Annexe D).

\subsubsection{Validation du modèle prédictif : Évaluation des temps de
réponse et précision (RMSE) IA vs
SUMO}\label{validation-du-moduxe8le-pruxe9dictif-uxe9valuation-des-temps-de-ruxe9ponse-et-pruxe9cision-rmse-ia-vs-sumo}

Pour contourner ces limites matérielles, le modèle prédictif basé sur
l'intelligence artificielle (XGBoost spectral) a été évalué en termes de
précision et de vitesse d'exécution face au simulateur physique SUMO de
référence.

\paragraph{Analyse comparative des temps de
calcul}\label{analyse-comparative-des-temps-de-calcul}

La différence de temps d'exécution entre l'approche physique
microscopique et l'inférence par apprentissage machine met en évidence
l'intérêt de la rupture technologique :

\begin{itemize}
\tightlist
\item
  \textbf{SUMO (Physique) :} Pour le réseau de Paris avec une charge de
  50 000 véhicules, l'exécution complète du pipeline (génération des
  routes via Dijkstra, calcul de la cinématique des agents et parsing
  final du fichier d'émissions XML \texttt{tripinfo.xml}) requiert un
  temps moyen de \textbf{45 minutes} de calcul CPU continu.
\item
  \textbf{Modèle Prédictif IA (XGBoost) :} L'inférence à partir des
  caractéristiques spectrales du graphe routier s'exécute en seulement
  \textbf{0,2 seconde} sur le même matériel, offrant un gain de vitesse
  de calcul significatif.
\end{itemize}

\paragraph{\texorpdfstring{Précision des estimations de pollution
(\(CO_2\))}{Précision des estimations de pollution (CO\_2)}}\label{pruxe9cision-des-estimations-de-pollution-co_2}

La précision du modèle prédictif a été quantifiée en calculant l'erreur
quadratique moyenne (Root Mean Squared Error - RMSE) sur un ensemble de
test indépendant composé de 200 scénarios urbains simulés sous
différentes configurations de volumes et de flottes :
\[RMSE = \sqrt{\frac{1}{N_{test}} \sum_{i=1}^{N_{test}} \left( y_{real, i} - y_{pred, i} \right)^2}\]

Les validations empiriques indiquent que le modèle XGBoost basé sur les
caractéristiques spectrales (rayon spectral, constante de Kreiss, etc.)
parvient à prédire le volume total d'émissions de \(CO_2\) avec un écart
moyen inférieur à \textbf{4,5 \%} par rapport aux résultats calculés par
le moteur physique de SUMO.

Cette précision démontre que la structure géométrique et spectrale du
graphe routier permet de capter la dynamique des flux sans avoir à
simuler individuellement chaque agent physique. Le planificateur urbain
dispose ainsi d'un outil d'évaluation instantané de l'empreinte
environnementale des aménagements projetés, combinant la réactivité de
l'IA et la fidélité de la micro-simulation.

\newpage

\section{CONCLUSION GÉNÉRALE ET
PERSPECTIVES}\label{conclusion-guxe9nuxe9rale-et-perspectives}

\subsubsection{Bilan des contributions
scientifiques}\label{bilan-des-contributions-scientifiques}

Ce travail de recherche a permis de jeter les bases d'un double cadre
méthodologique pour la modélisation de la mobilité urbaine décarbonée.

D'une part, nous avons développé des jumeaux numériques microscopiques
de haute-fidélité capables de modéliser l'impact cinématique local de
l'intégration d'infrastructures de recharge pour véhicules électriques.
Ce travail a été appliqué sur l'axe historique Paris-Versailles ainsi
que sur le quartier résidentiel à haute densité de Sao Bien à Hanoï. Il
a permis d'intégrer des comportements stochastiques de charge et de
documenter des dynamiques de trafic spécifiques comme l'inversion de
vacances (\emph{Holiday Reversal}).

D'autre part, nous avons proposé une méthode de rupture basée sur
l'intelligence artificielle topologique spectrale. En exploitant la
structure mathématique de la matrice d'adjacence du réseau urbain (rayon
spectral, constante de Kreiss, perturbations de Kato) et l'algorithme
XGBoost, nous avons démontré qu'il est possible d'estimer instantanément
la pollution d'une ville sans simuler individuellement chaque véhicule,
contournant ainsi les limites computationnelles traditionnelles.

\subsubsection{La boucle d'optimisation hybride (IA-SUMO) comme
perspective
ultime}\label{la-boucle-doptimisation-hybride-ia-sumo-comme-perspective-ultime}

La complémentarité des deux approches développées dans ce mémoire ouvre
la voie à un cadre d'optimisation urbaine hybride (IA-SUMO) alliant la
rapidité de l'apprentissage machine et la précision chirurgicale de la
simulation physique.

Dans ce schéma opérationnel : 1. Le modèle prédictif basé sur l'IA
(XGBoost) est utilisé en amont pour explorer rapidement de larges
espaces de solutions (par exemple, tester des milliers d'implantations
géométriques de voirie ou de localisations de hubs de recharge). L'IA
évalue chaque configuration en une fraction de seconde, éliminant les
scénarios inefficaces et sélectionnant les configurations optimales. 2.
Les scénarios retenus par le modèle prédictif sont ensuite injectés dans
le jumeau numérique microscopique haute-fidélité (SUMO). Cette
simulation physique détaillée permet d'affiner l'analyse locale
(vérification des files d'attente au mètre près, comportement
d'évitement des deux-roues, impact électrique précis).

Cette boucle hybride permet d'optimiser l'aménagement urbain à grande
échelle tout en conservant une précision physique sur les points
sensibles du réseau.

\subsubsection{Valorisation académique et publications
associées}\label{valorisation-acaduxe9mique-et-publications-associuxe9es}

Les résultats de ce travail de recherche font l'objet d'une valorisation
académique à travers la préparation de deux publications scientifiques
de fin d'année :

\begin{enumerate}
\def\labelenumi{\arabic{enumi}.}
\tightlist
\item
  \textbf{Première Publication (co-écrite avec VinUniversity) :}

  \begin{itemize}
  \tightlist
  \item
    \emph{Titre :} ``Microscopic Traffic Flow and Emission Modeling of
    High-Power Electric Vehicle Charging Infrastructure in Hyper-Dense
    Master-Planned Communities: The Case of Sao Bien, Vinhomes Ocean
    Park.''
  \item
    \emph{Contenu :} Présentation du protocole de capture par vision par
    ordinateur (YOLOv8), de l'architecture du jumeau numérique SUMO
    localisé, et de l'évaluation des scénarios de densification à Sao
    Bien.
  \end{itemize}
\item
  \textbf{Deuxième Publication (co-écrite avec l'École Hexagone) :}

  \begin{itemize}
  \tightlist
  \item
    \emph{Titre :} ``Topological Graph-Spectral Machine Learning for
    Real-Time Urban CO2 Emissions Prediction: Applying Kato's
    Perturbation Theory and Kreiss Constants to Non-Normal Urban
    Networks.''
  \item
    \emph{Contenu :} Formalisation de la méthode de prédiction spectrale
    sur graphes non-symétriques, étude de la stabilité via la constante
    de Kreiss, et évaluation des performances du modèle XGBoost par
    rapport au simulateur de référence.
  \end{itemize}
\end{enumerate}

\newpage

\section{GLOSSAIRE DES TERMES TECHNIQUES ET
MATHÉMATIQUES}\label{glossaire-des-termes-techniques-et-mathuxe9matiques}

\begin{itemize}
\tightlist
\item
  \textbf{Jumeau Numérique (Digital Twin) :} Réplication virtuelle
  dynamique d'un système physique réel (ici, la voirie et la cinématique
  des flux de trafic d'une ville) permettant d'effectuer des tests, de
  simuler des scénarios d'aménagements et de guider la prise de
  décision.
\item
  \textbf{Micro-simulation microscopique :} Modélisation individuelle du
  comportement de chaque agent mobile (vitesse, position, distance de
  sécurité) à chaque pas de temps discret, par opposition aux modèles
  macroscopiques basés sur des équations d'écoulement de fluides moyens.
\item
  \textbf{Mésoscopique :} Modèle de simulation hybride intermédiaire
  dans lequel les flux de trafic ne sont pas calculés véhicule par
  véhicule, mais modélisés sous forme de files d'attente de paquets
  d'agents se déplaçant le long de segments de voies, réduisant
  significativement la charge computationnelle.
\item
  \textbf{TraCI (Traffic Control Interface) :} API native du framework
  SUMO permettant le contrôle et la modification temps réel des états de
  simulation (feux, départs, itinéraires) via un script externe
  (généralement écrit en Python).
\item
  \textbf{Matrice d'Adjacence :} Représentation matricielle carrée de
  taille \(n \times n\) (où \(n\) est le nombre de nœuds) décrivant les
  connexions d'un graphe. Dans le cadre de l'ingénierie routière, elle
  est asymétrique pour modéliser les sens uniques et pondérée
  physiquement par le rapport temps de parcours/capacité des voies.
\item
  \textbf{Opérateur Laplacien :} Matrice définie par \(L = D - A\) (ou
  ses variantes normalisées) décrivant les flux de diffusion sur le
  graphe de la voirie. Ses valeurs propres caractérisent la connectivité
  structurelle du réseau.
\item
  \textbf{Rayon Spectral (\(\rho\)) :} Module de la valeur propre
  dominante d'une matrice. En topologie urbaine, il caractérise la
  capacité globale de transit et la hiérarchisation des corridors de
  circulation.
\item
  \textbf{Non-normalité :} Propriété d'une matrice qui ne commute pas
  avec sa transposée (\(A A^T \neq A^T A\)). Caractérise les réseaux
  urbains asymétriques et indique la possibilité d'amplifications
  transitoires du trafic.
\item
  \textbf{Constante de Kreiss (\(K\)) :} Métrique caractérisant
  l'amplitude maximale de l'amplification transitoire des perturbations
  avant le retour à l'équilibre asymptotique, traduisant la
  vulnérabilité d'un réseau aux embouteillages en cascade.
\item
  \textbf{Normes de Hardy (\(H_2 / H_\infty\)) :} Normes d'espaces
  fonctionnels décrivant la réponse d'un système dynamique aux
  perturbations. \(H_\infty\) caractérise le gain maximal dans le pire
  scénario de charge, tandis que \(H_2\) caractérise la persistance
  temporelle de l'énergie de perturbation stockée (mémoire de la
  congestion).
\item
  \textbf{Théorie des Perturbations de Kato :} Formalisme mathématique
  permettant de dériver analytiquement les variations du spectre d'un
  opérateur linéaire soumis à une perturbation mineure (par exemple,
  calculer l'impact de la fermeture d'un pont sur les valeurs propres de
  transit sans recalculer le système complet).
\item
  \textbf{XGBoost (eXtreme Gradient Boosting) :} Algorithme
  d'apprentissage supervisé basé sur le boosting d'arbres de décision
  régularisé, optimisant une fonction de coût complexe via un
  développement de Taylor de second ordre.
\item
  \textbf{YOLOv8 (You Only Look Once, v8) :} Modèle de réseau de
  neurones convolutif monocoup optimisé pour le traitement d'images
  temps réel, réalisant la détection, la classification et le suivi
  d'objets (véhicules).
\item
  \textbf{Dwell Time (Temps d'arrêt) :} Durée d'immobilisation physique
  d'un véhicule à une borne pour réaliser sa session de recharge
  électrique, modélisée ici sous forme de loi normale tronquée.
\item
  \textbf{Warm-Start (Démarrage à chaud) :} Initialisation d'une
  simulation dans un état pré-chargé (hub de recharge occupé
  stochastiquement par des véhicules fantômes) pour supprimer le biais
  transitoire du démarrage à vide.
\item
  \textbf{Spillback (Refoulement) :} Phénomène de propagation d'un
  bouchon où la file d'attente accumulée sur une voie sature et déborde
  pour paralyser les intesections situées en amont.
\item
  \textbf{SWAP (Pagination) :} Mécanisme du noyau système consistant à
  utiliser une partie de l'espace disque comme mémoire virtuelle lente
  lorsque la mémoire vive (RAM) physique est saturée.
\end{itemize}

\newpage

\section{BIBLIOGRAPHIE SCIENTIFIQUE DE
RÉFÉRENCE}\label{bibliographie-scientifique-de-ruxe9fuxe9rence}

\begin{itemize}
\tightlist
\item
  Agarwal, M., Maze, T. H., \& Souleyrette, R. (2005). \emph{Impacts of
  Weather on Urban Freeway Traffic Flow Characteristics and Facility
  Capacity}. Center for Transportation Research and Education, Iowa
  State University.
\item
  Aimsun (2024). \emph{Aimsun Next User Manual}. Aimsun SLU.
\item
  Apur (2025). \emph{Évolution des mobilités et du parc automobile dans
  la Métropole du Grand Paris au 1er Janvier 2025}. Atelier Parisien
  d'Urbanisme.
\item
  Avere-France (2025). \emph{Baromètre national des infrastructures de
  recharge de véhicules électriques}. Association nationale pour le
  développement de la mobilité électrique.
\item
  Boeing, G. (2017). \emph{OSMnx: New methods for acquiring,
  constructing, analyzing, and visualizing complex street networks}.
  Computers, Environment and Urban Systems, 65, 126-139.
\item
  Braess, D. (1968). \emph{Über ein Paradoxon aus der Verkehrsplanung}.
  Unternehmensforschung, 12, 258-268.
\item
  City of Paris (2024). \emph{Rapport d'activité sur les déplacements et
  la circulation à Paris en 2023}. Mairie de Paris.
\item
  Easley, D., \& Kleinberg, J. (2010). \emph{Networks, Crowds, and
  Markets: Reasoning about a Highly Connected World}. Cambridge
  University Press.
\item
  European Environment Agency (EEA) (2024). \emph{Transport and
  Mobility: Transitioning to Sustainable Systems}. EEA Report.
\item
  Forrester, J. W. (1969). \emph{Urban Dynamics}. Productivity Press.
\item
  French Government (2026). \emph{Bonus écologique et primes pour
  l'acquisition d'un véhicule électrique}. Service-Public.fr.
\item
  Green SM (2025). \emph{GSM: Green and Smart Mobility electric fleet
  expansion}. GSM Joint Stock Company.
\item
  International Energy Agency (IEA) (2025). \emph{Road Transport --
  Breakthrough Agenda Report 2025}. IEA.
\item
  Kato, T. (1995). \emph{Perturbation Theory for Linear Operators} (Vol.
  132). Springer Science \& Business Media.
\item
  Krajzewicz, D., Erdmann, J., Behrisch, M., \& Bieker, L. (2012).
  \emph{Recent development and applications of SUMO - Simulation of
  Urban MObility}. International Journal on Advances in Systems and
  Measurements, 5(3\&4), 128-138.
\item
  Kratzsch, V. et al.~(2020). \emph{HBEFA Version 4.1 -- Handbook
  Emission Factors for Road Transport}. INFRAS.
\item
  Kreiss, H. O. (1962). \emph{Über die Stabilitätsdefinition für
  Differenzengleichungen, die partielle Differentialgleichungen
  approximieren}. BIT Numerical Mathematics, 2(3), 153-181.
\item
  MATsim (2024). \emph{Multi-Agent Transport Simulation (MATSim)}.
  MATsim Community.
\item
  Mitchell, M. (2009). \emph{Complexity: A Guided Tour}. Oxford
  University Press.
\item
  OECD (2024). \emph{Urban Mobility and Decarbonisation}. Organisation
  for Economic Co-operation and Development.
\item
  PTV Group (2024). \emph{PTV Vissim User Manual: Microscopic Traffic
  Flow Simulation}. PTV AG.
\item
  Streamlit Inc.~(2020). \emph{Streamlit: The fastest way to build and
  share data applications for machine learning}.
\item
  SUMO (2024). \emph{SUMO User Documentation}. German Aerospace Center
  (DLR).
\item
  Texas A\&M Transportation Institute (2021). \emph{Urban Mobility
  Report}. Texas A\&M University.
\item
  Thocle (2025). \emph{Realistic Urban Traffic Generator using
  Decentralized Federated Learning for the SUMO simulator}.
  arXiv:2506.07980.
\item
  Trefethen, L. N., \& Embree, M. (2005). \emph{Spectra and
  Pseudospectra: The Behavior of Nonnormal Matrices and Operators}.
  Princeton University Press.
\item
  United Nations, Department of Economic and Social Affairs (2018).
  \emph{World Urbanization Prospects: The 2018 Revision}. United
  Nations.
\item
  United Nations, Department of Economic and Social Affairs (2018).
  \emph{Sustainable Urbanization is Key to Successful Development}.
  United Nations.
\item
  VinBus (2021). \emph{VinBus officially operates the first smart
  electric bus in Vietnam}. VinBus Ecology Transport Services.
\item
  Vinhomes (2024). \emph{Vinhomes Ocean Park: Urban Development and
  Master-Planned Communities}. Vinhomes Joint Stock Company.
\item
  World Bank (2024). \emph{Urban Development Overview}. World Bank
  Group.
\end{itemize}

\newpage

\section{ANNEXES : RÉSULTATS DÉTAILLÉS DES RUNS DE SIMULATION
MULTI-VILLES}\label{annexes-ruxe9sultats-duxe9tailluxe9s-des-runs-de-simulation-multi-villes}

Cette annexe compile l'ensemble des résultats numériques issus des 17
simulations complètes exécutées dans notre framework multi-villes
agnostique (simulation\_runs\_summary.md).

\subsubsection{Tableau des indicateurs macroscopiques des simulations
complètes}\label{tableau-des-indicateurs-macroscopiques-des-simulations-compluxe8tes}

{\def\LTcaptype{none} % do not increment counter
\begin{longtable}[]{@{}
  >{\raggedright\arraybackslash}p{(\linewidth - 14\tabcolsep) * \real{0.1053}}
  >{\raggedright\arraybackslash}p{(\linewidth - 14\tabcolsep) * \real{0.1053}}
  >{\centering\arraybackslash}p{(\linewidth - 14\tabcolsep) * \real{0.1316}}
  >{\centering\arraybackslash}p{(\linewidth - 14\tabcolsep) * \real{0.1316}}
  >{\centering\arraybackslash}p{(\linewidth - 14\tabcolsep) * \real{0.1316}}
  >{\centering\arraybackslash}p{(\linewidth - 14\tabcolsep) * \real{0.1316}}
  >{\centering\arraybackslash}p{(\linewidth - 14\tabcolsep) * \real{0.1316}}
  >{\centering\arraybackslash}p{(\linewidth - 14\tabcolsep) * \real{0.1316}}@{}}
\toprule\noalign{}
\begin{minipage}[b]{\linewidth}\raggedright
Date \& Time
\end{minipage} & \begin{minipage}[b]{\linewidth}\raggedright
City
\end{minipage} & \begin{minipage}[b]{\linewidth}\centering
Vehicles
\end{minipage} & \begin{minipage}[b]{\linewidth}\centering
Avg Speed
\end{minipage} & \begin{minipage}[b]{\linewidth}\centering
CO₂ Emitted
\end{minipage} & \begin{minipage}[b]{\linewidth}\centering
Weather
\end{minipage} & \begin{minipage}[b]{\linewidth}\centering
EV Adoption
\end{minipage} & \begin{minipage}[b]{\linewidth}\centering
Execution Time
\end{minipage} \\
\midrule\noalign{}
\endhead
\bottomrule\noalign{}
\endlastfoot
2026-06-08 11:37:40 & \textbf{Sydney} & 47 239 & 2.76 km/h & 174.23 Tons
& Clear (Optimal) & 15\% & 2602.90 s \\
2026-06-08 10:15:27 & \textbf{Rio de Janeiro} & 64 172 & 2.72 km/h &
265.52 Tons & Clear (Optimal) & 15\% & 3787.60 s \\
2026-06-08 09:58:55 & \textbf{Monaco} & 6 080 & 5.86 km/h & 21.46 Tons &
Clear (Optimal) & 8\% & 185.89 s \\
2026-06-06 10:55:18 & \textbf{Casablanca} & 49 785 & 15.62 km/h & 172.64
Tons & Clear (Optimal) & 10\% & 17605.75 s \\
2026-06-06 09:30:40 & \textbf{London} & 68 239 & 1.52 km/h & 241.86 Tons
& Clear (Optimal) & 15\% & 4545.16 s \\
2026-05-22 11:31:03 & \textbf{Sydney} & 55 453 & 2.21 km/h & 215.45 Tons
& Clear (Optimal) & 15\% & 3752.45 s \\
2026-05-22 10:33:58 & \textbf{Chamonix} & 2 000 & 46.50 km/h & 2.32 Tons
& Clear (Optimal) & 8\% & 13.07 s \\
2026-05-21 11:54:12 & \textbf{Versailles} & 1 000 & 31.46 km/h & 0.76
Tons & Clear (Optimal) & 10\% & 18.96 s \\
2026-05-13 20:22:36 & \textbf{Versailles} & 20 450 & 2.31 km/h & 82.06
Tons & Clear (Optimal) & 10\% & 521.54 s \\
2026-05-12 15:36:46 & \textbf{Madrid} & 98 560 & 16.19 km/h & 352.83
Tons & Clear (Optimal) & 10\% & 12276.98 s \\
2026-05-12 11:03:53 & \textbf{Paris} & 128 644 & 3.08 km/h & 508.20 Tons
& Clear (Optimal) & 10\% & 5971.15 s \\
2026-04-13 06:07:41 & \textbf{Paris} & 92 173 & 4.75 km/h & 338.01 Tons
& Clear (Optimal) & 10\% & 4126.56 s \\
2026-04-12 12:14:49 & \textbf{Paris} & 90 909 & 4.06 km/h & 341.46 Tons
& Clear (Optimal) & 10\% & 4340.81 s \\
2026-04-12 08:06:36 & \textbf{Paris} & 48 554 & 8.74 km/h & 152.21 Tons
& Clear (Optimal) & N/A & 2036.43 s \\
2026-04-11 16:47:14 & \textbf{Paris} & 48 641 & 7.93 km/h & 151.69 Tons
& Clear (Optimal) & N/A & 2051.00 s \\
2026-04-11 16:11:53 & \textbf{Versailles} & 12 750 & 4.26 km/h & 49.45
Tons & Clear (Optimal) & N/A & 219.36 s \\
2026-04-11 13:25:05 & \textbf{Paris} & 48 669 & 9.09 km/h & 166.46 Tons
& Clear (Optimal) & N/A & 2039.20 s \\
\end{longtable}
}

\subsubsection{\texorpdfstring{Tonnage de \(CO_2\) émis par classe de
véhicule de la
flotte}{Tonnage de CO\_2 émis par classe de véhicule de la flotte}}\label{tonnage-de-co_2-uxe9mis-par-classe-de-vuxe9hicule-de-la-flotte}

{\def\LTcaptype{none} % do not increment counter
\begin{longtable}[]{@{}
  >{\raggedright\arraybackslash}p{(\linewidth - 12\tabcolsep) * \real{0.1212}}
  >{\raggedright\arraybackslash}p{(\linewidth - 12\tabcolsep) * \real{0.1212}}
  >{\centering\arraybackslash}p{(\linewidth - 12\tabcolsep) * \real{0.1515}}
  >{\centering\arraybackslash}p{(\linewidth - 12\tabcolsep) * \real{0.1515}}
  >{\centering\arraybackslash}p{(\linewidth - 12\tabcolsep) * \real{0.1515}}
  >{\centering\arraybackslash}p{(\linewidth - 12\tabcolsep) * \real{0.1515}}
  >{\centering\arraybackslash}p{(\linewidth - 12\tabcolsep) * \real{0.1515}}@{}}
\toprule\noalign{}
\begin{minipage}[b]{\linewidth}\raggedright
Date \& Time
\end{minipage} & \begin{minipage}[b]{\linewidth}\raggedright
City
\end{minipage} & \begin{minipage}[b]{\linewidth}\centering
Total CO₂
\end{minipage} & \begin{minipage}[b]{\linewidth}\centering
Cars CO₂
\end{minipage} & \begin{minipage}[b]{\linewidth}\centering
Motorcycles CO₂
\end{minipage} & \begin{minipage}[b]{\linewidth}\centering
Buses CO₂
\end{minipage} & \begin{minipage}[b]{\linewidth}\centering
Trucks CO₂
\end{minipage} \\
\midrule\noalign{}
\endhead
\bottomrule\noalign{}
\endlastfoot
2026-06-08 11:37:40 & \textbf{Sydney} & \textbf{174.23 T} & 88.43 T &
37.37 T & 23.26 T & 25.18 T \\
2026-06-08 10:15:27 & \textbf{Rio de Janeiro} & \textbf{265.52 T} &
134.77 T & 56.41 T & 36.30 T & 38.04 T \\
2026-06-08 09:58:55 & \textbf{Monaco} & \textbf{21.46 T} & 17.24 T &
2.26 T & 0.95 T & 1.01 T \\
2026-06-06 10:55:18 & \textbf{Casablanca} & \textbf{172.64 T} & 104.59 T
& 26.36 T & 24.60 T & 17.09 T \\
2026-06-06 09:30:40 & \textbf{London} & \textbf{241.86 T} & 123.13 T &
53.54 T & 31.05 T & 34.14 T \\
2026-05-22 11:31:03 & \textbf{Sydney} & \textbf{215.45 T} & 109.64 T &
46.80 T & 28.13 T & 30.89 T \\
2026-05-22 10:33:58 & \textbf{Chamonix} & \textbf{2.32 T} & 1.88 T &
0.23 T & 0.11 T & 0.10 T \\
2026-05-21 11:54:12 & \textbf{Versailles} & \textbf{0.76 T} & 0.54 T &
0.12 T & 0.08 T & 0.03 T \\
2026-05-13 20:22:36 & \textbf{Versailles} & \textbf{82.06 T} & 57.07 T &
14.03 T & 7.12 T & 3.83 T \\
2026-05-12 15:36:46 & \textbf{Madrid} & \textbf{352.83 T} & 178.47 T &
79.62 T & 55.68 T & 39.05 T \\
2026-05-12 11:03:53 & \textbf{Paris} & \textbf{508.20 T} & 374.95 T &
105.01 T & 9.15 T & 19.10 T \\
2026-04-13 06:07:41 & \textbf{Paris} & \textbf{338.01 T} & 270.84 T &
47.96 T & 6.28 T & 12.93 T \\
2026-04-12 12:14:49 & \textbf{Paris} & \textbf{341.46 T} & 155.06 T &
88.45 T & 63.24 T & 34.71 T \\
2026-04-12 08:06:36 & \textbf{Paris} & \textbf{152.21 T} & 69.28 T &
38.10 T & 29.29 T & 15.53 T \\
2026-04-11 16:47:14 & \textbf{Paris} & \textbf{151.69 T} & 68.97 T &
38.21 T & 29.13 T & 15.38 T \\
2026-04-11 16:11:53 & \textbf{Versailles} & \textbf{49.45 T} & 0.00 T &
0.00 T & 0.00 T & 0.00 T \\
2026-04-11 13:25:05 & \textbf{Paris} & \textbf{166.46 T} & 0.00 T & 0.00
T & 0.00 T & 0.00 T \\
\end{longtable}
}

\subsubsection{Profil de performance informatique de l'exécution
(secondes
CPU)}\label{profil-de-performance-informatique-de-lexuxe9cution-secondes-cpu}

{\def\LTcaptype{none} % do not increment counter
\begin{longtable}[]{@{}
  >{\raggedright\arraybackslash}p{(\linewidth - 10\tabcolsep) * \real{0.1429}}
  >{\raggedright\arraybackslash}p{(\linewidth - 10\tabcolsep) * \real{0.1429}}
  >{\centering\arraybackslash}p{(\linewidth - 10\tabcolsep) * \real{0.1786}}
  >{\centering\arraybackslash}p{(\linewidth - 10\tabcolsep) * \real{0.1786}}
  >{\centering\arraybackslash}p{(\linewidth - 10\tabcolsep) * \real{0.1786}}
  >{\centering\arraybackslash}p{(\linewidth - 10\tabcolsep) * \real{0.1786}}@{}}
\toprule\noalign{}
\begin{minipage}[b]{\linewidth}\raggedright
Date \& Time
\end{minipage} & \begin{minipage}[b]{\linewidth}\raggedright
City
\end{minipage} & \begin{minipage}[b]{\linewidth}\centering
Dijkstra Routing
\end{minipage} & \begin{minipage}[b]{\linewidth}\centering
SUMO Physics Engine
\end{minipage} & \begin{minipage}[b]{\linewidth}\centering
XML Parsing
\end{minipage} & \begin{minipage}[b]{\linewidth}\centering
Total Duration
\end{minipage} \\
\midrule\noalign{}
\endhead
\bottomrule\noalign{}
\endlastfoot
2026-06-08 11:37:40 & \textbf{Sydney} & 666.64 s & 1934.63 s & 1.58 s &
\textbf{2602.90 s} \\
2026-06-08 10:15:27 & \textbf{Rio de Janeiro} & 555.46 s & 3229.83 s &
2.30 s & \textbf{3787.60 s} \\
2026-06-08 09:58:55 & \textbf{Monaco} & 7.69 s & 178.03 s & 0.16 s &
\textbf{185.89 s} \\
2026-06-06 10:55:18 & \textbf{Casablanca} & 13577.62 s & 4026.18 s &
1.89 s & \textbf{17605.75 s} \\
2026-06-06 09:30:40 & \textbf{London} & 1886.52 s & 2655.74 s & 2.90 s &
\textbf{4545.16 s} \\
2026-05-22 11:31:03 & \textbf{Sydney} & 997.87 s & 2752.49 s & 2.08 s &
\textbf{3752.45 s} \\
2026-05-22 10:33:58 & \textbf{Chamonix} & 4.27 s & 8.75 s & 0.04 s &
\textbf{13.07 s} \\
2026-05-21 11:54:12 & \textbf{Versailles} & 6.07 s & 12.84 s & 0.03 s &
\textbf{18.96 s} \\
2026-05-13 20:22:36 & \textbf{Versailles} & 48.35 s & 472.63 s & 0.55 s
& \textbf{521.54 s} \\
2026-05-12 15:36:46 & \textbf{Madrid} & 9716.06 s & 2558.17 s & 2.72 s &
\textbf{12276.98 s} \\
2026-05-12 11:03:53 & \textbf{Paris} & 3371.01 s & 2596.00 s & 4.13 s &
\textbf{5971.15 s} \\
2026-04-13 06:07:41 & \textbf{Paris} & 2353.06 s & 1770.98 s & 2.50 s &
\textbf{4126.56 s} \\
2026-04-12 12:14:49 & \textbf{Paris} & 2421.37 s & 1916.28 s & 3.16 s &
\textbf{4340.81 s} \\
2026-04-12 08:06:36 & \textbf{Paris} & 1167.02 s & 868.03 s & 1.37 s &
\textbf{2036.43 s} \\
2026-04-11 16:47:14 & \textbf{Paris} & 1134.40 s & 915.39 s & 1.20 s &
\textbf{2051.00 s} \\
2026-04-11 16:11:53 & \textbf{Versailles} & 17.36 s & 201.73 s & 0.26 s
& \textbf{219.36 s} \\
2026-04-11 13:25:05 & \textbf{Paris} & 1145.37 s & 892.65 s & 1.17 s &
\textbf{2039.20 s} \\
\end{longtable}
}

\subsubsection{Index des illustrations graphiques et des visualisations
de
simulation}\label{index-des-illustrations-graphiques-et-des-visualisations-de-simulation}

Cette section rassemble les planches graphiques du jumeau numérique, de
la topologie spectrale et du modèle d'apprentissage supervisé.

\paragraph{Planches relatives au hub de recharge de Sao Bien
(Hanoï)}\label{planches-relatives-au-hub-de-recharge-de-sao-bien-hanouxef}

\begin{itemize}
\tightlist
\item
  \textbf{Figure D.1 : Détection et classification multi-classes via
  YOLOv8}\\
  {[}Illustration de la vue plongeante du 3ème étage du Vincom Mega Mall
  montrant les boîtes de délimitation de détection pour les classes
  \emph{Standard Car}, \emph{Electric Vehicle} et \emph{Motorcycle} avec
  l'affichage de l'identifiant unique de suivi (tracking){]}.
\item
  \textbf{Figure D.2 : Géométrie d'accès du Hub de recharge de Sao Bien
  (Floor 0)}\\
  {[}Plan masse géométrique représentant les 12 places de recharge en
  épi à 135 degrés le long de la voie de service et la file d'attente de
  taxis GSM en amont{]}.
\item
  \textbf{Figure D.3 : Profils de charge de trafic et inversion modale
  (« Holiday Reversal »)}\\
  {[}Histogramme comparatif de la répartition modale de la flotte de
  véhicules légers en circulation à midi : Semaine classique (dominée
  par 50.7 \% de motos) vs Vacances nationales (dominée par 40.0 \% de
  véhicules électriques et 66.3 \% de part électrique sur le total des
  véhicules à 4 roues){]}.
\item
  \textbf{Figure D.4 : Profils d'évolution temporelle des temps
  d'attente au hub sous les scénarios S0 à S5}\\
  {[}Graphiques de files d'attente stochastiques montrant l'envolée des
  temps d'attente à 45.2 minutes sous le scénario critique de saturation
  S4b, et sa chute spectaculaire à 6.8 minutes sous le scénario de
  mitigation S5 par expansion physique (+8 slots){]}.
\end{itemize}

\paragraph{Planches de topologie spectrale et mathématique des réseaux
urbains}\label{planches-de-topologie-spectrale-et-mathuxe9matique-des-ruxe9seaux-urbains}

\begin{itemize}
\tightlist
\item
  \textbf{Figure D.5 : Distribution des spectres des matrices
  d'adjacence dans le plan complexe}\\
  {[}Nuages de points représentant les valeurs propres complexes
  (\(\lambda_1, \dots, \lambda_n\)) des graphes asymétriques pour les
  villes de Paris, Los Angeles et Hanoï, illustrant la dissymétrie et
  les structures communautaires{]}.
\item
  \textbf{Figure D.6 : Frontières de pseudospectres et non-normalité}\\
  {[}Contours de la résolvante \(\|(zI - A)^{-1}\|_2\) dans le plan
  complexe pour Marseille, Paris et Los Angeles, mettant en évidence les
  zones d'amplification transitoire et le calcul graphique de la
  constante de Kreiss{]}.
\end{itemize}

\paragraph{Planches d'évaluation du modèle d'Intelligence
Artificielle}\label{planches-duxe9valuation-du-moduxe8le-dintelligence-artificielle}

\begin{itemize}
\tightlist
\item
  \textbf{Figure D.7 : Importance relative des 43 descripteurs du modèle
  XGBoost (Feature Importance)}\\
  {[}Graphique en barres horizontales affichant les coefficients
  d'importance du modèle final : domination de \texttt{nb\_total\_veh}
  et \texttt{nb\_voitures}, suivi des métriques spectrales pondérées
  \texttt{sigma\_max\_weighted} et \texttt{h2\_norm\_weighted}{]}.
\item
  \textbf{Figure D.8 : Courbes d'apprentissage et de convergence de
  XGBoost}\\
  {[}Graphique de l'évolution de la fonction de perte (RMSE) sur les
  sous-ensembles d'entraînement et de validation au cours des itérations
  d'arbres (boosting rounds), démontrant la stabilité de la convergence
  sans sur-apprentissage{]}.
\item
  \textbf{Figure D.9 : Diagramme de dispersion (Scatter Plot) IA vs
  Ground Truth (SUMO)}\\
  {[}Graphique de dispersion comparant les émissions de CO₂ prédites par
  le modèle XGBoost face aux valeurs réelles mesurées dans le simulateur
  de référence SUMO, avec la ligne de régression \(y=x\) et le calcul de
  l'erreur moyenne de 4.5 \%{]}.
\end{itemize}

\paragraph{Planches de cartographie des congestions et heatmaps de
simulation}\label{planches-de-cartographie-des-congestions-et-heatmaps-de-simulation}

\begin{itemize}
\tightlist
\item
  \textbf{Figure D.10 : Cartes thermiques (Heatmaps) de congestion
  spatiale dans le jumeau numérique}\\
  {[}Visualisations SIG colorées des réseaux routiers de Paris, Londres,
  Los Angeles et Hanoï après une simulation de pointe de 1 heure. Les
  segments routiers sont colorés du gris foncé (fluide) au rouge
  écarlate (saturation complète / gridlock){]}.
\item
  \textbf{Figure D.11 : Comparaison des durées informatiques et de
  l'occupation mémoire RAM}\\
  {[}Diagramme à barres comparant l'occupation de la RAM (de 110 Mo pour
  Versailles à 18.5 Go pour Hanoï) et le temps de calcul CPU (45 minutes
  pour SUMO physique à Paris contre 0.2 seconde pour l'inférence
  XGBoost){]}.
\end{itemize}

\end{document}
